This document is normative for the semantics, and it names no implementation of one. Every claim
below stands on the argument written beside it; nothing here is true because something elsewhere
checks it.

\textbf{It states what is true now, and carries no history.}

\begin{center}\rule{0.5\linewidth}{0.5pt}\end{center}

\section{§1 The attribute domain}\label{1-the-attribute-domain}

A store holds items under names, and every value it holds, whatever type a schema gives it, is
drawn from one shared domain of attributes. This section fixes that domain.

Everything a table holds lives in one set. Types are read as retracts of it, which is the device
this whole specification rests on.

\begin{verbatim}
-- members, and the attribute domain over them
m  ∈ Member  ::= S | N | B | BOOL | NULL | SS | NS | BS | L | M
a  ∈ 𝔸       ::= ⟨m, v⟩        where v ∈ Carries(m)

-- an item is an M without its wrapper: the one place a map is not written as one
i  ∈ Item    ::= Name ⇀fin 𝔸
\end{verbatim}

The coproduct is \textbf{exact}: an attribute carries one member and no other. That is a fact about the
wire format and it is load-bearing here — \texttt{⟨S,\ "1"⟩} and \texttt{⟨N,\ "1"⟩} are distinct elements of \texttt{𝔸}, so
a codec choosing between them is choosing between distinguishable things.

\begin{tabular}{@{}lll@{}}
\toprule
Member & \texttt{Carries(m)} & \\
\midrule
\texttt{S} & \texttt{String} & not normalised \\
\texttt{N} & \texttt{Numerals} & optional sign, digits, optional point, optional exponent; ≤ 38 significant digits; magnitude 0 or \texttt{1e-130\ …\ 9.9…e125} \\
\texttt{B} & \texttt{Bytes} & \\
\texttt{BOOL} & \texttt{Bool} & \\
\texttt{NULL} & \texttt{𝟙} & present, and carrying nothing \\
\texttt{SS} & \texttt{℘⁺fin(String)} & non-empty; unordered \\
\texttt{NS} & \texttt{℘⁺fin(Numerals)} & non-empty; unordered \\
\texttt{BS} & \texttt{℘⁺fin(Bytes)} & non-empty; unordered \\
\texttt{L} & \texttt{𝔸*} & ordered; heterogeneous \\
\texttt{M} & \texttt{Name\ ⇀fin\ 𝔸} & \\
\bottomrule
\end{tabular}

\texttt{Numerals} admits an exponent — DynamoDB states its own range \cite{dynamodb-types} as \texttt{1E-130} to \texttt{9.9999…E+125}, so
\texttt{"1e+21"} is a well-formed \texttt{N} and not a defect. The bound that bites is \emph{magnitude}, not notation.

Note where \texttt{℘} appears and where it does not. \texttt{SS} is a \emph{set} and \texttt{L} is a \emph{sequence}; nothing in the
domain identifies them. Every claim later about telling a set from an array traces back to this line.

\begin{center}\rule{0.5\linewidth}{0.5pt}\end{center}

\section{§2 What a codec is}\label{2-what-a-codec-is}

A partial retraction of the attribute domain onto a set of values, together with the equivalence its
own reading induces.

\begin{verbatim}
c = (T, μ, e, p)      T ⊆ Value     μ ∈ Member ∪ {⊤}
                      e : T ⇀ 𝔸    encode — may fail
                      p : 𝔸 ⇀ T    decode — may fail
\end{verbatim}

\begin{verbatim}
v ∈ Value ::= str(String) | num(D) | bool(Bool) | bin(Bytes) | null
            | strSet(℘⁺fin(String)) | numSet(℘⁺fin(D)) | binSet(℘⁺fin(Bytes))
            | arr(Value*) | map(Name ⇀fin Value)

k ∈ Kind  ::= string | number | boolean | binary | set | array | map | null
\end{verbatim}

\texttt{T} is a set of values, the language of things a codec holds. Ten constructors, and eight
kinds: the three set constructors are all of kind \texttt{set}, and every other constructor is a kind
of its own. \texttt{D} is the doubles, fixed in §3. \texttt{Value} is not §1's \texttt{Member}: an \texttt{N} carries
\texttt{Numerals}, which is text, where \texttt{num} carries a double, and §3's \texttt{decimal} is the term that
closes the gap. \texttt{μ\ =\ ⊤} reads as
\emph{not determined by the term}, not as \emph{anything}. Which terms may leave it
undetermined is settled in §3.

\subsection{Laws}\label{laws}

\begin{definition}[(R) retraction]
\label{cond:R}
\lean{Calculus.Lawful.retract}
\leanok
\texttt{∀\ t\ ∈\ dom(e).\ p(e(t))\ =\ t}. Writing then reading returns the value written.
\end{definition}

\begin{definition}[(C) canonizability]
\label{cond:C}
\lean{Calculus.Lawful.canon}
\leanok
\texttt{∀\ a\ ∈\ dom(p).\ p(a)\ ∈\ dom(e)}. Whatever was read can be written back.
\end{definition}

\begin{definition}[(W) a term writes its member]
\label{cond:W}
\lean{Calculus.Lawful.writes}
\leanok
\texttt{μ\ ≠\ ⊤\ ⟹\ ∀\ t\ ∈\ dom(e).\ e(t)\ =\ ⟨μ,\ \_⟩}. A term that names its member writes that member
and no other.
\end{definition}

\begin{definition}[(MP) member preservation]
\label{cond:MP}
\uses{cond:C}
\lean{Calculus.Lawful.agree, Calculus.Agree.member_eq}
\leanok
\texttt{∀\ a\ ∈\ dom(p).\ member(e(p(a)))\ =\ member(a)}. Reading and writing back never changes a stored
member. §5's strong form \texttt{e(p(a))\ ≼\ a} implies this.
\end{definition}

\begin{theorem}[the four laws are independent]
\label{thm:laws-independent}
\uses{cond:R,cond:C,cond:W,cond:MP}
\lean{Calculus.cR_notR, Calculus.cR_C, Calculus.cR_W, Calculus.cR_MP, Calculus.cC_R, Calculus.cC_notC, Calculus.cC_W, Calculus.cC_MP, Calculus.cW_notW, Calculus.cW_R, Calculus.cW_C, Calculus.cW_MP, Calculus.cM_R, Calculus.cM_C, Calculus.cM_W, Calculus.cM_notMP}
\leanok
No one of \rulelink{cond:R}{\textbf{(R)}}, \rulelink{cond:C}{\textbf{(C)}}, \rulelink{cond:W}{\textbf{(W)}} and \rulelink{cond:MP}{\textbf{(MP)}} follows from the other three: for each
there is a pair satisfying the other three and failing that one.
\end{theorem}

And \rulelink{cond:MP}{\textbf{(MP)}} is well defined only because \rulelink{cond:C}{\textbf{(C)}} holds — \texttt{e(p(a))} is an
attribute at all because everything readable re-encodes. One term is exempt from \rulelink{cond:MP}{\textbf{(MP)}} and only
one: \texttt{legacy}, whose entire purpose is to change what is stored (§3, §6), and since \rulelink{cond:MP}{\textbf{(MP)}} is
recursive, a term \emph{containing} a \texttt{legacy} inherits the exemption on exactly the attributes the
migration rereads, and on no others. Nothing is exempt from the other three.

Only \rulelink{cond:R}{\textbf{(R)}} is an equality on the nose. Taken as one, \texttt{e\ ∘\ p\ =\ id} is \emph{false} and must be: digits have
a normal form, so writing back what was read is free to normalise it, and a schema may discard names
it does not declare. In the language of quotient lenses \cite{quotient-lenses} \texttt{e\ ∘\ p} is a
\textbf{canonizer} — total on \texttt{dom(p)} by \rulelink{cond:C}{\textbf{(C)}} — and the equivalence it works modulo is not a parameter
a programmer chooses: it is \texttt{ker\ p}, forced by the reading itself. Write \texttt{a\ ≈c\ a′\ ⟺\ p(a)\ =\ p(a′)}.

\subsection{Stability under a generator}\label{corollary-q--stability}

\begin{corollary}[Q — stability]
\label{cor:Q}
\lean{Calculus.Lawful.stable}
\leanok
\texttt{∀\ a\ ∈\ dom(p).\ e(p(a))\ ≈c\ a}.
\end{corollary}

\begin{proof}
\uses{cond:C,cond:R}
\leanok
By \rulelink{cond:C}{\textbf{(C)}}, \texttt{p(a)\ ∈\ dom(e)}. Put \texttt{t\ =\ p(a)}. By \rulelink{cond:R}{\textbf{(R)}}, \texttt{p(e(t))\ =\ t}. Substituting,
\texttt{p(e(p(a)))\ =\ p(a)}, which is \texttt{e(p(a))\ ≈c\ a} by definition of \texttt{≈c}.
\end{proof}

So stability is retraction asked at a different \emph{source of values}: \rulelink{cond:R}{\textbf{(R)}} at values a generator
produced, \rulelink{cor:Q}{\textbf{(Q)}} at values a table returned. What \rulelink{cor:Q}{\textbf{(Q)}} adds to
\rulelink{cond:R}{\textbf{(R)}} is exactly \rulelink{cond:C}{\textbf{(C)}}, and that containment is the thing most likely to be untrue.

\subsection{What the readable side owes}\label{corollary-read-side-w}

\begin{corollary}[read-side W]
\label{cor:readW}
\lean{Calculus.Lawful.reads}
\leanok
\texttt{μ\ ≠\ ⊤\ ⟹\ ∀\ a\ ∈\ dom(p).\ member(a)\ =\ μ}.
\end{corollary}

\begin{proof}
\uses{cond:C,cond:W,cond:MP}
\leanok
\texttt{p(a)\ ∈\ dom(e)} by \rulelink{cond:C}{\textbf{(C)}}, so \texttt{e(p(a))\ =\ ⟨μ,\ \_⟩} by \rulelink{cond:W}{\textbf{(W)}}, and
\texttt{member(a)\ =\ member(e(p(a)))\ =\ μ} by \rulelink{cond:MP}{\textbf{(MP)}}.
\end{proof}

\rulelink{cond:W}{\textbf{(W)}} constrains only the writing half; this corollary is what makes the reading half follow. A
\texttt{⇝\ S} term does not merely write \texttt{S}: it \emph{accepts} nothing else, or a read-modify-write could
smuggle a member change through. \texttt{legacy} is precisely the sanctioned violation: it reads members the
current encoding no longer writes, and re-encoding them is the migration.

\rulelink{cond:MP}{\textbf{(MP)}} binds hardest where \texttt{μ\ =\ ⊤}. A term with a \texttt{μ} is pinned by \rulelink{cond:W}{\textbf{(W)}} and the corollary above;
a term without one chooses per value, and \rulelink{cond:R}{\textbf{(R)}} and \rulelink{cor:Q}{\textbf{(Q)}} do not stop it choosing differently on
the way back. A codec reading \texttt{⟨N,\ "1.5"⟩} to a string and writing \texttt{⟨S,\ "1.5"⟩} satisfies both — the
\emph{value} survives — while silently changing the stored type, so a numeric sort key stops ordering and a
typed index rejects the write. It holds per attribute and recursively: a member change three levels
inside an \texttt{M} is the same failure.

This is why \texttt{opaque} refuses \texttt{N} unless \texttt{ν} admits it. Its refusal is \rulelink{cond:MP}{\textbf{(MP)}}, and §2 without it could not say
so.

Write \texttt{Can(c)\ =\ e{[}T{]}} for the \textbf{canonical} attributes: those a codec will actually produce. \texttt{e\ ∘\ p}
is then the canonizer, mapping any readable attribute to the canonical representative of its class. A
codec is total on \texttt{Can(c)} and merely partial on \texttt{𝔸}, which is the precise sense in which a table is
more permissive than the schema describing it.

\subsection{What a codec ideally is}\label{what-a-codec-ideally-is}

Those pieces assemble into something simpler than the four laws suggest, and that is where the
complexity lives.

\textbf{Between \texttt{dom(e)} and \texttt{Can(c)}, a codec is a bijection.} \texttt{p\ ∘\ e\ =\ id} there is \rulelink{cond:R}{\textbf{(R)}}; \texttt{e\ ∘\ p\ =\ id}
there follows from it, since a canonical attribute is \texttt{e(t)} for some \texttt{t} and \rulelink{cond:R}{\textbf{(R)}} sends it back
to that \texttt{t}; and \texttt{e} is injective for the same reason. No quotient, no partiality, nothing to
choose. On that part a codec is a plain isomorphism between a set of values and a set of attributes,
and there is nothing further to say about it.

Everything else in §2 is one of exactly \textbf{two extensions} of that bijection, in the two directions
it can be extended:

\begin{tabular}{@{}llll@{}}
\toprule
& extension & what it buys & the law that governs it \\
\midrule
\textbf{quotient} & \texttt{dom(p)\ ⊇\ Can(c)}, strictly & reading back what the store hands you, not only what you wrote & \rulelink{cond:C}{\textbf{(C)}}, and \texttt{≈c} \\
\textbf{partiality} & \texttt{dom(e)\ ⊆\ T}, \texttt{dom(p)\ ⊊\ 𝔸} & refusing values and attributes that are not yours & the domains themselves \\
\bottomrule
\end{tabular}

So \texttt{e\ ∘\ p\ =\ id} being false is not a defect of the pair, and not a weakening of the iso. It is the
first extension, seen from outside: on \texttt{Can(c)} the equation \emph{does} hold, and off it the canonizer
maps each class to its representative. The literature has a name for each — quotient lenses \cite{quotient-lenses} for
the first, partial isomorphisms \cite{invertible-syntax} for the second — and this calculus needs both, which is why
neither reference alone is the right one. The quotient reading refines the asymmetric lens of
\cite{tree-transformations}, and \cite{lens-laws} is where the variants are told apart.

The same two relaxations reappear one level up, in the terms rather than in the pair: §3's \texttt{held} is
an iso and \texttt{refine} its prism, and §6's \texttt{legacy(c,\ r)\ ≈\ c} on \texttt{Can(c)} is the quotient again at the
term level. That is not a coincidence but it is not a theorem either: the algebra that builds
codecs relaxes an iso in the two ways a codec itself does.

\subsection{When the deviations vanish}\label{when-the-deviations-vanish}

Both extensions can be driven to zero, and asking what it takes is the cleanest way to see what
each one is really for. They do not have the same answer, and in this calculus they cannot be
answered at once.

\textbf{The quotient collapses} when \texttt{ker\ p} is equality — when no two distinct attributes read to one
value. That fails for exactly three reasons, and one of them is about the medium:

\begin{itemize}
\tightlist
\item
  \textbf{\texttt{η}, the store's own normalisation of a stored numeral (§3), normalises \texttt{N}}. This is the medium, and it is the calculus's one concession to
  something it does not model.
\item
  \textbf{\texttt{drop} discards undeclared names} (§4). This is a schema's own decision, taken deliberately.
\item
  \textbf{A set loses a payload} when two of its payloads read to one value (§5). This is the
  container's own, and \rulelink{cond:R}{\textbf{(R)}} is untouched, because the canonical payloads are the ones the
  codec wrote.
\end{itemize}

So a codec collapses the quotient by holding no set, reading no normalised numeral, and declaring
every name it means to keep. \texttt{decimal} is the witness that this is reachable: its held set is \texttt{Stored}, the numerals
\texttt{η} leaves alone (§3), and its \texttt{ker\ p} is
equality, which is precisely why its held set had to be narrowed to \texttt{Stored}.

\textbf{Partiality collapses} when \texttt{dom(p)\ =\ 𝔸} — when every attribute the domain admits is one this
codec reads. And here the answer is structural rather than a matter of care: \textbf{§1 fixes one \texttt{𝔸}
for many \texttt{T}.} A table holds items of several schemas; a codec is a retract of the domain, not a
description of it. So partiality is not a tax on textual representation, and no choice of format
removes it. It is the price of the domain being \emph{shared}, and it would be paid identically by a
free algebra, a fixed-width binary layout, or a memory image, as long as more than one \texttt{T} lives in
one \texttt{𝔸}. What removes it is one schema per store, which is a different architecture rather than a
better encoding.

\textbf{They do not collapse together.} \texttt{decimal} kills the quotient and keeps partiality; \texttt{opaque(lossy)}
comes closest to reading all of \texttt{𝔸} and keeps the quotient, since its \texttt{N} option reads through the
numeral-to-double reading. No term of this calculus has both at once, and §1 is why.

One deviation \emph{does} have a clean term-level answer, and it is often mistaken for the other
two. Ambiguity in a \texttt{union} is not a property of the medium at
all — it is what happens when a choice is written without a tag. \rulelink{rule:dunion}{\texttt{(dunion)}} is the tagged case, and
its \rulelink{cond:U}{\textbf{(U)}} is neither an obligation — a condition the semantics enforces per
input by refusing (§5) — nor a theorem: the tag reaches one variant by
construction. Whatever the storage looks like, that
collapse is available by declaring the discriminator.

\begin{center}\rule{0.5\linewidth}{0.5pt}\end{center}

\section{§3 The judgement}\label{3-the-judgement}

Two indices, because a codec has two sides: the values it holds, and the member it stores them in.
Both are needed, and neither determines the other.

\begin{verbatim}
⊢ c : T ⇝ μ        c holds T, and stores it as μ
\end{verbatim}

\subsection{Ground terms}\label{ground-terms}

\begin{definition}[(string)]
\label{rule:string}
\lean{Calculus.Typed.string}
\leanok
\begin{verbatim}
──────────────────────────────
⊢ string : String ⇝ S      (string)
\end{verbatim}
\end{definition}

\begin{definition}[(decimal)]
\label{rule:decimal}
\lean{Calculus.Typed.decimal}
\leanok
\begin{verbatim}
──────────────────────────────
⊢ decimal : Stored ⇝ N     (decimal)
\end{verbatim}
\end{definition}

\begin{definition}[(number)]
\label{rule:number}
\lean{Calculus.Typed.number}
\leanok
\begin{verbatim}
D = { x ∈ Double : x = 0 ∨ 1e-130 ≤ |x| ≤ 9.9…e125 }
────────────────────────────────────────────────────── (number)
⊢ number : D ⇝ N
\end{verbatim}
\end{definition}

\begin{definition}[(boolean)]
\label{rule:boolean}
\lean{Calculus.Typed.boolean}
\leanok
\begin{verbatim}
──────────────────────────────
⊢ boolean : Bool ⇝ BOOL    (boolean)
\end{verbatim}
\end{definition}

\begin{definition}[(binary)]
\label{rule:binary}
\lean{Calculus.Typed.binary}
\leanok
\begin{verbatim}
──────────────────────────────
⊢ binary : Bytes ⇝ B       (binary)
\end{verbatim}
\end{definition}

\begin{definition}[(null)]
\label{rule:null}
\lean{Calculus.Typed.null}
\leanok
\begin{verbatim}
──────────────────────────────
⊢ null : {null} ⇝ NULL     (null)
\end{verbatim}
\end{definition}

\begin{definition}[(literal)]
\label{rule:literal}
\lean{Calculus.Typed.literal}
\leanok
\begin{verbatim}
⊢ g : T ⇝ μ,  g ground    ℓ₁ … ℓₙ ∈ T,  n ≥ 1,  pairwise distinct
───────────────────────────────────────────────────────────────── (literal)
⊢ literal(ℓ₁ … ℓₙ) : {ℓ₁ … ℓₙ} ⇝ μ
\end{verbatim}
\end{definition}

\texttt{D} is a set of doubles; the numeric host is taken as given and not modelled here.
\texttt{number} is \textbf{not} restricted to integers, and not to \texttt{2⁵³}. Every finite double satisfies \rulelink{cond:R}{\textbf{(R)}},
because \texttt{String} on a double is by definition the shortest text reading back as that same double — so
\texttt{1.5} round-trips, and so does \texttt{2⁵³\ +\ 1}, which \emph{is} \texttt{2⁵³} and round-trips as that. What bounds \texttt{D} is
DynamoDB's magnitude range: at most 38 significant digits, and magnitude zero or \texttt{1e-130\ …\ 9.9…e125}.

\texttt{D} is \textbf{checked}, on both rules that write an \texttt{N} from a double. A term that writes an attribute
outside §1's \texttt{𝔸} is not a term of this calculus, so the bound cannot be a precondition recorded and
left to the service.

The two rules check it in the two different places their admissible sets live. \rulelink{rule:number}{\texttt{(number)}} gates
\textbf{both halves}: \texttt{e} on the text the double is written as, and \texttt{p} on the double the text reads
back to, which is not redundant — a numeral the store holds can read back as a double past the
maximum, and gating only the write is the configuration that breaks \rulelink{cond:C}{\textbf{(C)}}.

Both-halves gating is not an extra duty these rules carry; it is what a gate \emph{is} in this calculus.
A predicate checked on the way in and on the way out is \texttt{refine}, and both number rules are \texttt{refine}
— \texttt{number} at \texttt{D}, \texttt{decimal} at \texttt{Stored}. Neither base is an \emph{ungated} one, and neither can be: with
the \texttt{N} carrier typed, an ungated numeric codec would have to write an attribute the domain has no
point for. \texttt{number} is itself the base on the double side, because §1's \texttt{Numerals} is the whole of
\texttt{Carries(N)}; \texttt{decimal}'s base is the numerals gated at well-formedness, which is the largest
decimal codec the carrier permits. Nothing about them is bespoke, and
the shape is the argument: an ad-hoc \texttt{if} can gate one half and a \texttt{refine} cannot. \rulelink{rule:literal}{\texttt{(literal)}} gates
at the \textbf{declaration}, alongside \texttt{n\ ≥\ 1} and pairwise distinctness, because \texttt{ℓ} is data the term
owns and a premise on data is decidable when the codec is built.

That premise discharges the \textbf{writable} side and only the writable side: \texttt{ℓ} is fixed when the codec is
built, so every value \texttt{e} will ever accept has been checked already. It discharges nothing on the
read side. Reading an \texttt{N} as a value is wider than membership in \texttt{ℓ}, so an attribute that is not a
numeral at all can land on a value the set contains, and the check passes on something §1 has no
point for. \textbf{A premise bounds the values a term owns, and a stored attribute is not one of them.}
So a numeric \texttt{literal} reads through the same gate as \rulelink{rule:number}{\texttt{(number)}}, and gets it from the same place:
its base \textbf{is} \texttt{number}, and \texttt{ℓ\ ⊆\ T} is the rule's premise, so the stored text must be a numeral
before it is read as a value.

A subset of \texttt{D} chosen by the caller — a finite set of admitted numbers — is \texttt{refine(number,\ T)}, and has no
rule of its own because it needs none. A term whose \texttt{p} returned the double it read \emph{without}
verifying membership in \texttt{T} would break \rulelink{cond:C}{\textbf{(C)}}, and with it the claim that \texttt{:\ T} says anything; a
term that does verify is \texttt{refine(number,\ T)} on all four fields. An unchecked narrowing of a read is
a cast, and a cast belongs to \texttt{refine}, which does
the check, or to \texttt{held}, which is where the caller's own functions live.

\texttt{decimal} holds numerals, not a machine number. That is the one ground term whose held set is the
stored text rather than the machine number, and §1's table is the reason: \texttt{N} carries 38
significant digits and a double carries 15.

Its held set is \texttt{Stored}, and not §1's \texttt{Numerals}, because those are \textbf{different sets}. \texttt{Numerals}
is what the store \emph{accepts}; \texttt{Stored} is what it \emph{hands back}. DynamoDB normalises every \texttt{N} it
takes — \texttt{1.500} becomes \texttt{1.5}, \texttt{007} becomes \texttt{7}, \texttt{1e5} becomes \texttt{100000}, \texttt{-0} becomes \texttt{0} — and
\texttt{decimal}'s \texttt{p} is the identity on the digits, so anything outside the normal form is written as one
string and read as another.

This is the calculus's one concession to something it does not model. §2 composes \texttt{e} and \texttt{p}
directly; in use there is a store between them, and it is not the identity on attributes. Writing
\texttt{η} for the store's normalisation, a codec survives the round trip through it exactly when \textbf{\texttt{η}
lies inside \texttt{ker\ p}} — \texttt{≈c} of §2, the relation that already exists for "a different attribute
reading to the same value". \texttt{number} satisfies that for free, since \texttt{p(1.500)} and \texttt{p(1.5)} are the
same double; \texttt{decimal} cannot, because its \texttt{ker\ p} is equality. So the held set is narrowed to the
fixed points of \texttt{η}, which is what \texttt{Stored} is. \texttt{η} is observed, not stipulated: it is the store's
behaviour, and the calculus can only be arranged around it.

\subsection{A ground term either takes its member whole, or it is a refinement}\label{a-ground-term-either-takes-its-member-whole-or-it-is-a-refinement}

Read the ground terms together and they fall into two kinds, with nothing in between. \texttt{string},
\texttt{boolean}, \texttt{binary} and \texttt{null} take their member's carrier \textbf{whole} — every string is a \texttt{String},
every bool a \texttt{Bool} — and there is nothing for them to check. The other four want a \emph{subset} of what
the store will hold at that member, and each is \texttt{refine} over one of those bases at the predicate
that names the subset:

\begin{tabular}{@{}lll@{}}
\toprule
term & base & predicate \\
\midrule
\texttt{number} & §1's \texttt{Numerals} & \texttt{D} (§3) \\
\texttt{decimal} & numerals that are well-formed by §1's grammar & \texttt{Stored} (§3) \\
\texttt{literal(ℓ)} & the ground term \texttt{g} & \texttt{∈\ ℓ} \\
\bottomrule
\end{tabular}

This is not a resemblance. A predicate wanted on \textbf{both} halves \emph{is} \texttt{refine}, because \texttt{e} and \texttt{p}
are partial and refusing is what the error channel does; a rule that gates one half and not the
other is the shape §3 spends a paragraph refusing at \rulelink{rule:number}{\texttt{(number)}}. So a gate is never a rule's private
business, and its predicate belongs in the term's admissible set, which is why a \texttt{number}'s
held values are \texttt{D} and a \texttt{decimal}'s are \texttt{Stored}, beside a literal's \texttt{ℓ} (§5): gating both
ways is a term deciding its own membership.

\textbf{Containers are the other kind, and they are not refinements.} \texttt{set}, \texttt{array}, \texttt{tuple}, \texttt{record}
and \texttt{object} narrow by their \emph{subterms} and their \emph{names}, never by a predicate on the whole value,
so there is no \texttt{q} for an admissible set to carry. §5 records the same fact from the other side: a
container admits its empty case whatever its element admits, so no set of one container's own values
separates it from another. Scalars are refined by predicates; containers are refined structurally, and only the first kind
has a gate to record.

\texttt{literal} is the one ground term whose member follows from its \emph{values} rather than from its name,
and the rule says so by reading both off a base: \texttt{ℓ} must lie in some ground term's held set, and
what that term stores is what the literal stores. One rule covers \texttt{S} and \texttt{N}, and covers \texttt{BOOL},
\texttt{B} and \texttt{NULL} with nothing added — a boolean literal is \texttt{literal(true)}, at base \texttt{boolean}. The
base is an index the derivation carries and the term does not write. An alternative's kind determines
it, except between \texttt{string} and \texttt{decimal}, which both hold text: \texttt{decimal} is what to write when text
is wanted stored as \texttt{N}, and the two derivations of one written literal differ in nothing else.

The rule is also not doing anything \texttt{refine} does not, which is why \texttt{literal}
reads as the abbreviation it is. \texttt{literal(ℓ)} is \texttt{refine(g,\ ∈\ ℓ)}, and the general fact behind every
literal union is one statement rather than a rule per member:

\begin{quote}
For \textbf{any derivable base}, a union of refinements with pairwise disjoint predicates admits
uniquely on both sides.
\end{quote}

\texttt{literal} at \texttt{string} and at \texttt{number} are two instantiations; \texttt{boolean} and \texttt{binary} are two more;
and the base need not be ground at all — two disjoint refinements of one \texttt{array(string)} separate by
the same argument. \emph{One base with disjoint predicates} is what the rule turns on, not the
literal.
What \texttt{literal} adds over writing \texttt{refine} is that \texttt{ℓ} is data the term owns, so both admissible
sets carry it (§5) and a reader sees the enumeration.

There is deliberately no derivation for a mixed set: a codec writes one member, and one \texttt{g} cannot
hold both a string and a double, so \texttt{literal(ℓ,\ k)} with \texttt{ℓ\ ∈\ String} and \texttt{k\ ∈\ D} has none. That is \rulelink{cond:W}{\textbf{(W)}} and not an omission, and it is not
a limitation either — the spelling is \texttt{union(literal(ℓ),\ literal(k))}, whose options separate on
both sides for free, by distinct members on the readable one and disjoint kinds on the writable.
What the union concludes is \texttt{⊤}, which is the honest index for a term that writes two members. A union of
literals all at the number base is what gives the term for an enum indexed by number — the held set is the literal union and the
stored member is \texttt{N}, checked in both directions — and note that \texttt{⇝\ N} carries \texttt{Numerals} on the wire
while the held set is numeric, so \rulelink{cond:R}{\texttt{(R)}} obliges the reading to return the number and not its text.

\subsection{Containers}\label{containers}

\begin{definition}[(set)]
\label{rule:set}
\uses{rule:string,rule:number,rule:decimal,rule:binary}
\lean{Calculus.Typed.set}
\leanok
\begin{verbatim}
b ∈ {string, number, decimal, binary}
─────────────────────────────────────── (set)
⊢ set(b) : ℘⁺fin(b.held) ⇝ b.member
\end{verbatim}
\end{definition}

\rulelink{rule:set}{\texttt{(set)}}'s premise line is \textbf{an index}, not a hypothesis: the element's
member is read off the base rather than tested, so the four bases the rule admits are exactly those that carry one,
\texttt{b.held} and \texttt{b.member} are functions of the index, and an element that names no member is a
type error. \textbf{Four bases onto three wire
members}: \texttt{number} and \texttt{decimal} both store \texttt{NS} and are not the same base: one holds
machine doubles, the other the numeral text.

What a base rather than a term costs: \texttt{set(refine(b,\ q))} and \texttt{set(literal(…))} are \textbf{not} lost — a
restriction of the elements is a restriction of the set — but
\texttt{set(held(…))} and \texttt{set(legacy(…))} are, having no set-level spelling and no appearance in any
schema.

\begin{definition}[(array)]
\label{rule:array}
\lean{Calculus.Typed.array}
\leanok
\begin{verbatim}
⊢ c : T ⇝ μ
──────────────────────── (array)
⊢ array(c) : T* ⇝ L
\end{verbatim}
\end{definition}

\begin{definition}[(tuple)]
\label{rule:tuple}
\lean{Calculus.Typed.tuple}
\leanok
\begin{verbatim}
⊢ cᵢ : Tᵢ ⇝ μᵢ    i ∈ 1…n
───────────────────────────────────── (tuple)
⊢ tuple(c₁ … cₙ) : T₁ × … × Tₙ ⇝ L
\end{verbatim}
\end{definition}

\begin{definition}[(tuple-tail)]
\label{rule:tuple-tail}
\lean{Calculus.Typed.tuple}
\leanok
\begin{verbatim}
⊢ cᵢ : Tᵢ ⇝ μᵢ  (i ∈ 1…n)    ⊢ c : T ⇝ μ
───────────────────────────────────────────────────────── (tuple-tail)
⊢ tuple(c₁ … cₙ, tail(c)) : T₁ × … × Tₙ × T* ⇝ L
\end{verbatim}
\end{definition}

\begin{definition}[(record)]
\label{rule:record}
\lean{Calculus.Typed.record}
\leanok
\begin{verbatim}
⊢ c : T ⇝ μ
──────────────────────────────── (record)
⊢ record(c) : Name ⇀fin T ⇝ M
\end{verbatim}
\end{definition}

\texttt{tail} has no rule of its own, which is the point: it is not a term but a \emph{position marker}.
It appears only in the premise of \rulelink{rule:tuple-tail}{\texttt{(tuple-tail)}}, and once — so
\texttt{tuple(tail(c),\ c′)} and \texttt{tuple(tail(c),\ tail(c′))} have no derivation, and neither does
\texttt{array(tail(c))} nor a \texttt{tail} in any position §4 or §5 admits. A list is already every element; a map
has names and not positions.

The premise on \texttt{set} is the whole reason \texttt{μ} has to be an index and not an afterthought. A set's
member is \emph{computed} from its element's, so an element that does not name one names no set, and the
rule refuses it rather than a runtime check discovering it. Nothing else in the calculus reads its
subterm's member this way.

\subsection{Wrappers}\label{wrappers}

\begin{definition}[(held)]
\label{rule:held}
\lean{Calculus.Typed.held, Calculus.IsoCancels}
\begin{verbatim}
⊢ c : I ⇝ μ    f : T → I    g : I → T    g ∘ f = id_T    f ∘ g = id_I
─────────────────────────────────────────────────────────────────────── (held)
⊢ held(c, f, g) : T ⇝ μ
\end{verbatim}
\end{definition}

\begin{definition}[(refine)]
\label{rule:refine}
\lean{Calculus.Typed.refine}
\leanok
\begin{verbatim}
⊢ c : T ⇝ μ    q ⊆ T
──────────────────────── (refine)
⊢ refine(c, q) : q ⇝ μ
\end{verbatim}
\end{definition}

\begin{definition}[(legacy)]
\label{rule:legacy}
\lean{Calculus.Typed.legacy, Calculus.ReadsCoherently}
\leanok
\begin{verbatim}
⊢ c : T ⇝ μ    r : 𝔸 ⇀ T    ran(r) ⊆ dom(e)    a ∈ dom(r) ∩ dom(p) ⟹ r(a) = p(a)
──────────────────────────────────────────────────────────────────────────────────── (legacy)
⊢ legacy(c, r) : T ⇝ μ
\end{verbatim}
\end{definition}

\begin{definition}[(lazy)]
\label{rule:lazy}
\lean{Calculus.Typed.lazy}
\leanok
\begin{verbatim}
γ : 𝟙 → Codec    ⊢ γ(∗) : T ⇝ μ
─────────────────────────────── (lazy)
⊢ lazy(γ) : T ⇝ ⊤
\end{verbatim}
\end{definition}

\begin{definition}[(custom)]
\label{rule:custom}
\uses{cond:R,cond:C,cond:MP}
\lean{Calculus.Typed.custom, Calculus.Obligations}
\leanok
\begin{verbatim}
e : T ⇀ 𝔸    p : 𝔸 ⇀ T    (R), (C), (MP) hold of (e, p)
───────────────────────────────────────────────────────── (custom)
⊢ custom(e, p) : T ⇝ ⊤
\end{verbatim}
\end{definition}

\begin{definition}[(opaque)]
\label{rule:opaque}
\lean{Calculus.Typed.opaq}
\leanok
\begin{verbatim}
ν ∈ {refuse, lossy}
─────────────────────────── (opaque)
⊢ opaque(ν) : Uν ⇝ ⊤
\end{verbatim}
\end{definition}

These \texttt{⊤}s are forced for different reasons, and the distinction matters. \texttt{lazy} could be asked but answering would force the thunk, which is the recursion it exists to
break. \texttt{custom} is a pair of opaque functions and there is nothing to ask, which is also why its
premise is the laws themselves: nothing structural guarantees them, so the rule demands them.
Demanding \rulelink{cond:R}{\textbf{(R)}} alone would under-demand, deriving a \texttt{⊤} term free to change stored members — the
exact term §2's argument for \texttt{opaque} says cannot exist. (\rulelink{cond:W}{\textbf{(W)}} is vacuous at \texttt{⊤}, so the three
named are the whole obligation.)

\texttt{opaque} is the calculus's custom-class term, and it is primitive only in presentation: it
unfolds to the \textbf{union of everything self-identifying, closed under the containers} — a fixpoint of shape
\texttt{union(string,\ boolean,\ binary,\ ⟨NULL⟩,\ number?,\ set(opaque),\ array(opaque),\ record(opaque))} —
whose \rulelink{cond:U}{\textbf{(U)}} holds because the options' kinds and members are pairwise disjoint. \texttt{Uν} is \emph{defined}
as the values whose member is recoverable from the value alone, so \rulelink{cond:R}{\textbf{(R)}} holds by construction;
what \texttt{ν} decides is \texttt{N}, the one member a bare value cannot answer for, and the fixpoint reading
gives the policy its sharpest statement: \textbf{\texttt{refuse} is the absence of the \texttt{number} option.} Nothing
is guessed, because there is nothing present to guess with. Under \texttt{lossy} the option is present and
\texttt{number\ ⇄\ N}: the member survives every read-modify-write, and the digits survive only while they
fit a double.

\texttt{held} and \texttt{refine} are the same shape at different strengths: an \textbf{iso} and a \textbf{prism}. The
literature collapses them — invertible syntax descriptions \cite{invertible-syntax} builds everything from an algebra of
partial isomorphisms, which is the merged combinator.

They are one term, in the following sense. Write the merged combinator
as an adapter that may refuse on either half — \texttt{f} in front of the writing, \texttt{g} behind the reading:

\begin{verbatim}
e = λv. f(v) ⨾ e_c        p = λa. p_c(a) ⨾ g
\end{verbatim}

\texttt{held(c,\ f,\ g)} is that with both halves \textbf{total}, and \texttt{refine(c,\ q)} is that with a \textbf{partial
identity} on both halves — refuse off \texttt{q}, pass through on it. The pairs agree on the nose, not up
to anything. So \texttt{refine} is not a second kind of adapter; it is the one adapter with a weaker
premise, and its whole content is that refusing \emph{is} available, because \texttt{e} and \texttt{p} are partial to
begin with.

They are kept apart for a reason outside the calculus, and only that one: a total \texttt{g} lets a
generator for \texttt{I} be pushed forward to one for \texttt{T}, while a partial one leaves only rejection
sampling. Same term, different search. The split is a note about testing, written into the syntax
so it cannot be lost, which is the same bargain \texttt{nullable} strikes with \texttt{union(c,\ null)}, and
should be read the same way.

\texttt{legacy}'s \texttt{r} is a bare partial map and not a codec, because only one direction is ever used: the
rule writes through \texttt{c} whatever it read through. Asking for a codec would ask for a half that cannot
be called, and would \emph{force} one wherever the older encoding's held set differs. The two side
premises are the soundness proof's levies (§5). \texttt{ran(r)\ ⊆\ dom(e)} is \rulelink{cond:C}{\textbf{(C)}}'s: a migration must
land on values the current encoding can write, or the read succeeds and the write-back is the
failure. The \textbf{coherence} premise — \texttt{r} agrees with \texttt{p} wherever both are defined — is what makes
\texttt{legacy} a migration and not a disguise: without it the term is a union with a privileged write
path, resolved by a precedence between its two readings, which is precisely the resolution order
\rulelink{cond:U}{\textbf{(U)}} exists to refuse. With it, \texttt{p\ ∪\ r} is a plain union of partial maps, nothing is privileged,
and a migrated row behaves \emph{by law} like a current one at the surface. How the migration achieves
that is the migration's duty, but that it must is the calculus's.

\texttt{nullable} carries \textbf{no side premise}. The
danger is real — a subterm moving \texttt{null} onto another member, or another member onto \texttt{null}, would
let the wrapper's \texttt{NULL} branch shadow a real value, failing \rulelink{cond:R}{\textbf{(R)}} in one direction and \rulelink{cond:MP}{\textbf{(MP)}}
in the other, and the fix is not a premise but the wrapper's own definition. \texttt{nullable(c)} adds the
null point \textbf{unless \texttt{c} already owns it}, asked of both halves at once:

\begin{verbatim}
owns(c)  ⟺  null ∈ dom(e_c)  ∨  ⟨NULL, ∗⟩ ∈ dom(p_c)
\end{verbatim}

Where \texttt{owns(c)}, the pair is \texttt{c}'s untouched. Where it does not, \texttt{c} reaches neither end of the null
point, and then \textbf{null-honesty is a theorem}: if \texttt{e\_c} wrote
\texttt{⟨NULL,\ ∗⟩}, \rulelink{cond:R}{\textbf{(R)}} would make \texttt{p\_c} read it straight back, and \texttt{owns(c)} would hold; if \texttt{p\_c} read
some attribute to \texttt{null}, \rulelink{cond:C}{\textbf{(C)}} would make \texttt{e\_c} write the point, and \texttt{owns(c)} would hold again.
Both hypotheses are refuted, not required.

The guard must be \textbf{joint}. Asked per side, \texttt{c} could decline to write \texttt{null} while still reading
\texttt{⟨NULL,\ ∗⟩} — the null point's two ends would belong to different options, the round trip would go
out through one and back through the other, and \rulelink{cond:R}{\textbf{(R)}} would break with nothing looking ambiguous
on either side taken alone.

None of this is a new kind of condition. \texttt{nullable(c)} is structurally the two-option union
\texttt{union(c,\ ⟨NULL⟩-point)}, and its honesty conditions are that pair's §5 side conditions —
uniqueness of admission on the way in, on the way out — which the guard makes hold rather than
demand.

\rulelink{rule:null}{\texttt{(null)}} above makes \texttt{union(c,\ null)} a term, so both spellings exist, and the relation between
them is a \textbf{dichotomy}, not a trade-off. It is the wrapper's definition:

\begin{itemize}
\tightlist
\item
  where \texttt{c} does not touch the null point — \texttt{e\_c(null)} and \texttt{p\_c(⟨NULL,\ ∗⟩)} both undefined —
  \texttt{nullable(c)} and \texttt{union(c,\ null)} are the \textbf{same codec};
\item
  where \texttt{c} does touch it, \texttt{nullable(c)} is \textbf{equal to \texttt{c}}: the wrapper adds
  nothing. \texttt{nullable(opaque)\ =\ opaque}, \texttt{nullable(nullable(c))\ =\ nullable(c)}, which §6's
  idempotence node states as an equation, and \texttt{nullable(union(…,\ null))\ =\ union(…,\ null)}.
\end{itemize}

So \texttt{nullable} is either the union or the identity, and the precedence that distinguishes it buys
nothing a schema would want. (The two are also \texttt{≡} unconditionally; they
differ as codecs only in definedness, and only where the wrapper is redundant. The single term
outside both readings is \texttt{nullable(null)} — member \texttt{⊤} over the \texttt{NULL} point, of no use.)

Which settles what \texttt{nullable} is: \textbf{sugar} — idempotent insertion of the null option into a pair
that may already have it. It expresses nothing new, it costs no premise, and its idempotence is a
property of insertion rather than a fact about \texttt{𝟙}. The one term outside both branches of the
dichotomy is \texttt{nullable(null)}, which under the guard is simply \texttt{null}.

\texttt{union(c,\ null)} replaces it at a subterm whose shape is unknown — \texttt{held}, \texttt{custom},
\texttt{lazy}: no discharger — no decidable condition sufficient for \rulelink{cond:U}{\textbf{(U)}} at the
declaration (§5) — means an \textbf{obligation}, not a refusal, so \texttt{union(c,\ null)} is
derivable there and carries \rulelink{cond:U}{\textbf{(U)}} per input.

\begin{center}\rule{0.5\linewidth}{0.5pt}\end{center}

\section{§4 Fields are not codecs}\label{4-fields-are-not-codecs}

A second judgement, and the stratification is the point: presence is a property of a name binding,
never of an attribute.

\begin{verbatim}
⊢ f : T ⊗ π ⇝ μ    π ∈ {req, opt}     f holds T, with presence π, stored as μ
\end{verbatim}

A field carries a member like anything else — an optional attribute is still stored somewhere, and
an index key needs to know where. What a field carries that an attribute does not is
the \texttt{⊗\ π}, and that is the whole stratification: a rule demanding \texttt{⊢\ c\ :\ T\ ⇝\ μ} cannot be given
something whose judgement has a \texttt{⊗}.

\begin{definition}[(field)]
\label{rule:field}
\lean{Calculus.TypedField.req}
\leanok
\begin{verbatim}
⊢ c : T ⇝ μ
────────────────── (field)
⊢ c : T ⊗ req ⇝ μ
\end{verbatim}
\end{definition}

\begin{definition}[(optional)]
\label{rule:optional}
\lean{Calculus.TypedField.optional}
\leanok
\begin{verbatim}
⊢ c : T ⇝ μ
────────────────────────────── (optional)
⊢ optional(c) : T ⊗ opt ⇝ μ
\end{verbatim}
\end{definition}

\begin{definition}[(object)]
\label{rule:object}
\lean{Calculus.Typed.object}
\leanok
\begin{verbatim}
⊢ fᵢ : Tᵢ ⊗ πᵢ  (i ∈ 1…n)    n₁ … nₙ distinct    ⊢ ρ : rest
───────────────────────────────────────────────────────────────────────── (object)
⊢ object({n₁ : f₁ … nₙ : fₙ}, ρ) : ⟨n₁ :^π₁ T₁ … nₙ :^πₙ Tₙ⟩ ⊕ R(ρ) ⇝ M
\end{verbatim}
\end{definition}

\begin{definition}[(rest)]
\label{rule:rest}
\lean{Calculus.TypedRest}
\leanok
\begin{verbatim}
────────────────────────────────────────────────────────── (rest)
⊢ ρ : rest  ⟸  ρ = drop | ρ = refuse | ⊢ ρ : T ⇝ μ
\end{verbatim}
\end{definition}

Here \texttt{n\ :\^{}opt\ T} is a binding the record may omit, \texttt{n\ :\^{}req\ T} one it may not, and \texttt{R(ρ)} is the open
part, joined to the declared names by \texttt{⊕}, which adds bindings no declared name claims: a
declared name is always held by its own set, so \texttt{⊕} is total and there is no collision to resolve.
The open part is
\texttt{∅} for \texttt{drop} and \texttt{refuse}, and \texttt{Name\ ⇀fin\ T} for a rest codec holding \texttt{T}. \texttt{drop} and
\texttt{refuse} differ only in what \texttt{p} does with an undeclared name — drop it, or fail — which is a
difference in \texttt{dom(p)} and not in \texttt{T}. And \texttt{drop} is a \textbf{filter, never a traversal mode}: the
discard is decided by the name alone, before anything is decoded, so it cannot depend on a value and
the round trip through the declared names is untouched by what was discarded. That is the quotient
of §2 again — an item is \texttt{≈c} itself minus its undeclared names — and, \texttt{SS}/\texttt{NS}/\texttt{BS} carrying
no ordering to lose, one of only three places \texttt{ker\ p} is not equality. The others are \texttt{η} and a set.

\subsection{The stratification does work}\label{the-stratification-does-work}

No rule derives \texttt{⊢\ optional(c)\ :\ T\ ⇝\ μ}. Presence lives in \texttt{⊗\ π} and nowhere else, so \texttt{optional} is a field.
\texttt{array}'s premise is an attribute judgement. Nor \texttt{set}, \texttt{tuple}, \texttt{record}, nor a union option, for the
same reason: none of those positions has a name for anything to be absent from.

\begin{center}\rule{0.5\linewidth}{0.5pt}\end{center}

\section{§5 Choice, and its side conditions}\label{5-choice-and-its-side-conditions}

The rules a grammar cannot state, because they relate siblings rather than a term to its subterms,
and, before them, the one semantic condition even the judgement cannot decide, which everything else
in this section exists to approximate.

\subsection{Unique admission}\label{unique-admission}

A union's pair is defined by \textbf{unique admission}. Each option answers, per input and per side,
whether the input could be its — an attribute on the way out, a value on the way in — and the union
acts through the option that answers alone. Write \texttt{V(c)\ ⊆\ Value} for the values an option may be
asked to write and \texttt{A(c)\ ⊆\ 𝔸} for the attributes it may be asked to read.

\begin{definition}[(U) unique admission]
\label{cond:U}
\lean{Calculus.UAtCan, Calculus.UAtAct}
\leanok
\begin{verbatim}
strong      ∀ i. ∀ t ∈ dom(eᵢ).  {j : t ∈ V(cⱼ)} = {i}
                              ∧  {j : eᵢ(t) ∈ A(cⱼ)} = {i}

restricted  ∀ i. ∀ t ∈ dom(eᵢ).  {j : t ∈ V(cⱼ)} = {i}
                              ⟹  {j : eᵢ(t) ∈ A(cⱼ)} = {i}
            ∀ i. ∀ a ∈ dom(pᵢ).  {j : a ∈ A(cⱼ)} = {i}
                              ∧  pᵢ(a) ∈ dom(eᵢ)
                              ⟹  {j : pᵢ(a) ∈ V(cⱼ)} = {i}
\end{verbatim}
Resolving one side alone resolves the other, the same option. Asked where the union \emph{acts}: at
each value it writes, and each attribute it reads. Ambiguity is a failure, never a rank. The letter
names the strong form: at every canonical pair of every option, that option alone admits the
value, and that option alone admits the attribute. What \rulelink{rule:union}{\texttt{(union)}} carries as an
obligation is weaker, the same demand restricted to the inputs the union itself resolves, and the
strong form implies it.
\end{definition}

\begin{definition}[(D) incomparable held sets]
\label{cond:D}
\lean{Calculus.DIncomp}
\leanok
\texttt{∀\ i\ ≠\ j.\ (¬(Tᵢ\ ⊆\ Tⱼ)\ ∧\ ¬(Tⱼ\ ⊆\ Tᵢ))\ ∨\ Tᵢ\ ∩\ Tⱼ\ =\ ∅}.
Well-formedness: of any two options, neither held set subsumes the other, or the two share no
value. Required of every pair, disjointness rejects working unions, which is why it is offered as
an alternative and not imposed; two disjoint options are distinguishable in the strongest way
there is. An empty option is the unit of the sum and reaches the condition this way, being
disjoint from everything: it contributes nothing to the union and so cannot make it ambiguous,
while being a subset of every held set, which incomparability alone would reject. Decidable from
the held sets alone. It does \textbf{not} discharge the readable side, for the reason the
machinery of admission gives.
\end{definition}

\begin{definition}[(union)]
\label{rule:union}
\uses{cond:D}
\lean{Calculus.Typed.union}
\leanok
\begin{verbatim}
⊢ cᵢ : Tᵢ ⇝ μᵢ  (i ∈ 1…n)    (D)
─────────────────────────────────────────────────────────────────── (union)
⊢ union(c₁ … cₙ) : ⋃ Tᵢ ⇝ (μ if ∀i. μᵢ = μ;  else ⊤)
\end{verbatim}
\end{definition}

\begin{definition}[(dunion)]
\label{rule:dunion}
\uses{rule:object,rule:literal}
\begin{verbatim}
⊢ oᵢ = object(sᵢ, ρᵢ) : Tᵢ ⇝ M    sᵢ(d) = literal(tᵢ)    tᵢ pairwise distinct
─────────────────────────────────────────────────────────────────────────────── (dunion)
⊢ discriminatedUnion(d, o₁ … oₙ) : ⋃ Tᵢ ⇝ M
\end{verbatim}
\end{definition}

Note that \rulelink{rule:dunion}{\texttt{(dunion)}}'s premise asks \texttt{sᵢ(d)} to \emph{be} the literal — an attribute term, hence a \texttt{req}
field. A row without the discriminator does not select a core option. A caller that accepts such
an old row supplies one outer \texttt{legacy} reader for the whole union.

\textbf{Arity one at the discriminator}, and n-ary literals everywhere else. The rule asks for a single
tag \texttt{tᵢ} per variant, and disjointness of singletons is distinctness. Soundness does not need the
restriction: \rulelink{rule:literal}{\texttt{(literal)}} is n-ary, a set of tags is a perfectly good
admissible set, and pairwise disjointness is the condition (U) needs.

What required it is what a discriminated union \emph{is}. It is a \textbf{lookup}: the stored tag chooses
the variant. A variant answering to two tags is a key that does not identify its variant, and a
reader who has the tag still cannot say which shape they hold without consulting the table twice.

Read \rulelink{cond:U}{\textbf{(U)}} carefully, because the scope is the whole content. \texttt{∣\{i\ :\ cᵢ\ admits\ it\}∣\ =\ 1} at
\emph{every} input would be a stronger demand, and a wrong one: where two options both admit, the union
is simply undefined, and a law quantified over \texttt{dom} cannot reach an input the pair declines. What
the laws need is that ambiguity never appears \emph{on the far side of a round trip the union did take} —
\texttt{e(t)} unambiguous on the way out obliges \texttt{e(t)} to be unambiguous on the way back. An input the
union refuses outright asks nothing of anyone.

The difference decides real schemas, and the next example is one.
\texttt{union(array(literal("A")),\ array(literal("B")))} has both options admitting \texttt{{[}{]}}, so it declines to
write the empty array, but every value it does write is read back by the one option that wrote it.
Under the stronger reading the union is refused; under this one it is lawful, minus the empty case.

Declaration order is never consulted, which is what puts commutativity in §6, and why \texttt{legacy}'s
coherence premise (§3) matters beyond \texttt{legacy}: a precedence between readings anywhere in the
calculus would be a union resolved by rank, which is what \rulelink{cond:U}{\textbf{(U)}} refuses.

\rulelink{cond:U}{\textbf{(U)}} quantifies over inputs, so it is not decidable at the declaration, where a
codec is built and its subterms are known but no input is. The rest of
this section is its \textbf{dischargers} — decidable conditions, each sufficient for \rulelink{cond:U}{\textbf{(U)}} on the
unions it covers:

\begin{tabular}{@{}lll@{}}
\toprule
Discharger & Decides & From \\
\midrule
distinct members & the readable side, where the stored attributes differ & \texttt{μ}, the judgement's own index \\
disjoint option admission, \texttt{∀\ i\ ≠\ j.\ V(cᵢ)\ ∩\ V(cⱼ)\ =\ ∅} on the writable side and \texttt{A(cᵢ)\ ∩\ A(cⱼ)\ =\ ∅} on the readable one & either side, each side asked on its own & admission and occupancy (one coarse label per term per side), below \\
tag disjointness & \rulelink{rule:dunion}{\texttt{(dunion)}}, outright — both sides & the discriminator literals \\
disjoint predicates over one base & refinements of a common \texttt{c}, outright — both sides & the predicates \texttt{q} \\
kind disjointness & \texttt{opaque}'s fixpoint (§3), outright & one option per kind of §2 \\
\bottomrule
\end{tabular}

\textbf{Every row is one condition at a different projection.} Fix a partial map from inputs to labels —
a stored member, a runtime kind, a discriminator's tag — and give each option the set of labels it
answers to. Where those sets are pairwise disjoint, no input can reach two options, because the
label is a property of the \emph{input} and lands in at most one set. That is the whole argument.

Said more directly, without projections at all: \textbf{each option gets one set of values it may write
and one set of attributes it may read, and separation is disjointness of those two families.} A
projection is a way of \emph{building} such a set — a runtime shape gives the values whose shape it
allows, \texttt{μ} the attributes whose member it allows, a discriminator's tag the items carrying one of
its tags. \textbf{Option admission} is the family itself, the finest there is: the set the union
consults when it asks whether an option could take the input.

Read that way the admission row stops offering a \emph{choice} between a shape and a set. The two constraints
simply \textbf{intersect} — an option may write \texttt{v} only if its shape allows the constructor of \texttt{v} \emph{and}
it admits \texttt{v} — and disjointness of the intersections follows from disjointness in either factor. A
coarse shape and an exact set are two factors of one family rather than two mechanisms, which is why
one row carries both. The two families are also independent, which is the point the pairwise reading
obscures: the writable side may be separated by a shape while the readable side is separated by a
member, and nothing has to reconcile the two.

What is left to distinguish the rows is one thing only: whether a row constrains \emph{both} families or
one. \texttt{μ} labels attributes and nothing else, so the member row settles the readable side and never
the writable one. A discriminator's tag and §1's kind label both, which is why those two settle
\rulelink{cond:U}{\textbf{(U)}} outright. The admission row is asked once per side and may settle one side or both; a
union is resolved when both sides are resolved, by whatever combination of rows reaches them.

Read that way, \rulelink{rule:dunion}{\texttt{(dunion)}} is not a second kind of sum. It is \rulelink{rule:union}{\texttt{(union)}} with admission \textbf{read
off} the discriminator instead of attempted, and its \rulelink{cond:U}{\texttt{(U)}} needs no discharge, because reading a
label off an input reaches one variant by construction.

What no discharger covers stays an obligation the semantics enforces per input by failing — safe for
every law except \rulelink{cond:R}{\textbf{(R)}} at a canonical attribute, and that residue is what the theorem
below names. An option whose shape is \texttt{⊤} is one the admission row never separates, for the same
reason it cannot be checked: it could be anything. That costs nothing, because failing to discharge is not
refusing: \rulelink{cond:U}{\textbf{(U)}}
still guards each value the option is actually given.

\subsection{A weaker condition, and what it costs}\label{a-weaker-condition-proved-and-not-adopted}

The reading above rejects \texttt{∣\{i\ :\ cᵢ\ admits\ it\}∣\ =\ 1} as too strong and settles for \emph{undefined at the
overlap}. There is a third position between them.

\textbf{Confluent admission.} Where several options admit, require only that they \textbf{agree}, and then
act. Stated as \emph{first-match plus an obligation}, which is what makes it usable: the rule takes the
first admitter and never compares attributes, while the obligation says the choice did not matter.
So declaration order enters the definition and leaves again as a theorem: the first-match rule is
one half, the agreement obligation the other.

It is strictly weaker than \rulelink{cond:U}{\textbf{(U)}}, so strictly more total, and the case it gains is the one this
section already named as the price: \texttt{union(array(literal("A")),\ array(literal("B")))} writes \texttt{{[}{]}},
because both options write \texttt{⟨L,\ {[}{]}⟩} and the choice is inconsequential. "Lawful, minus the empty
case" loses the \emph{minus}.

What it costs is stated as plainly. \textbf{Order-independence stops being unconditional.} Under \rulelink{cond:U}{\textbf{(U)}}
commutativity is a fact about the definition: §6's \texttt{union(c₁,\ c₂,\ …)\ ≡\ union(c₂,\ c₁,\ …)} has no
hypotheses. Under confluence it holds \emph{given the obligation}, and a specification whose commutativity
rests on a premise is weaker as a specification, whatever it gains in totality. That trade is the
reason this is recorded here and not adopted.

And it changes what the fixpoint argument says. That a union whose options are genuinely recursive
has no fixpoint semantics is a refutation of \textbf{counting}, not of \texttt{union}:

\begin{theorem}[confluent admission survives the refutation]
\label{thm:confluent-survives}
\uses{cond:U}
\lean{Calculus.confluent_survives_the_refutation}
\leanok
In the witness both options write \texttt{⟨S,"A"⟩}, so the weaker rule still resolves where counting
collapses to \texttt{none}.
\end{theorem}

\begin{theorem}[confluent admission is monotone]
\label{thm:confluent-mono}
\uses{cond:U}
\lean{Calculus.selectF_mono}
\leanok
Monotone given the obligation, which counting provably lacks.
\end{theorem}

\begin{theorem}[confluent admission is continuous]
\label{thm:confluent-continuous}
\uses{cond:U}
\lean{Calculus.selectF_continuous_list, Calculus.selectF_continuous}
\leanok
Continuous, and in general this needs the obligation \textbf{at each stage} too. The two-option case
escapes it by putting the growing option first, which makes it automatically the first admitter.
\end{theorem}

Counting's continuity \emph{is} unconditional — its singleton requirement at the limit forces
uniqueness at every earlier stage — and counting cannot be monotone at all. So one requirement, the
obligation holding at every stage, buys both; what unique admission cannot buy at any price is
monotonicity, being a \emph{negative} condition on the other options.

Two things that are \textbf{not} established, and the second is why \texttt{lazy} under \texttt{union} is
outside the theorem. A fixpoint is not shown to exist: monotonicity and continuity are its obligations, not
its construction, and both theorems are the two-option case. And the obligation must hold at every
stage of an unrolling, which is not automatic: two \emph{ground} terms both take \texttt{.str\ "1.5"}
and store it differently. Where the options are pairwise disjoint it is
vacuous, so a syntactic disjointness test on the options is what would license it.

\textbf{Normatively, the rule is unique admission as stated above.} This subsection records a proved
alternative and the trade, not a change to the judgement.

\subsection{The machinery of admission}\label{the-machinery-of-admission}

The admission row asks, of each option and each side, which inputs the option may be given:
\texttt{V(c)} and \texttt{A(c)} above. Two things build them, both by structural recursion over the terms of §3–§4, and
both total.

\begin{theorem}[a term's admissible sets are exact]
\label{thm:admissible-exact}
\lean{Calculus.Adm_exact, Calculus.ExactV, Calculus.AdmA_exact, Calculus.ExactA}
\leanok
A codec is a pair, and an admissible set is a claim about one half of it at a time, so there are two
and not one: the values \texttt{e} accepts, and the attributes \texttt{p} accepts. Each is a \textbf{domain
rather than a description of one} — \texttt{dom(e)} and \texttt{dom(p)} — at every term of the calculus, with
no hypothesis and no term excepted. Nothing is over-approximated, and no term's answer is \texttt{⊤}.
\end{theorem}

That is available only because a row may \textbf{ask the codec} rather than describe it, and the four
terms force the point:

\begin{itemize}
\tightlist
\item
  \textbf{A numeric literal's readable side.} Distinct numerals read to the same double, so no set of
  values bounds a set of attributes. Nor do the stored texts, which \emph{are} known at the declaration:
  \texttt{literal(10²¹)} is written \texttt{⟨N,\ "1e+21"⟩} and handed back \texttt{⟨N,\ "1000000000000000000000"⟩}, which
  only the numeral-to-double read recovers, so a syntactic set would refuse a row the codec itself
  wrote — \texttt{ker\ p} shrinking below \texttt{η}. The set is \texttt{dom(p)} of the very pair the rule builds.
\item
  \textbf{\texttt{refine(c,\ q)}.} \texttt{q} is a predicate on values, so the writable side is \texttt{V(c)\ ∩\ q} directly. The
  readable side cannot follow syntactically, \texttt{q} testing the value a read \emph{produces}, so it decodes
  with the base and applies \texttt{q} to the result.
\item
  \textbf{\texttt{legacy(c,\ r)}.} \texttt{p} reads attributes \texttt{c} does not, so a passthrough would claim a bound on
  exactly the direction \texttt{r} widens: \texttt{legacy(literal("A"),\ r)} with \texttt{r} reading \texttt{⟨N,\ "1"⟩} to \texttt{"A"}
  reads an attribute \texttt{A(literal("A"))} refuses, which is unsound rather than coarse. The attributes
  \texttt{r} widens by are \texttt{dom(r)}, and \texttt{r} is data the term carries, so the set is \texttt{A(c)\ ∪\ dom(r)}. The
  writable side passes through: \texttt{legacy} adds no way to write.
\item
  \textbf{\texttt{held(c,\ f,\ g)}}, the one term whose readable side sees further than its writable one. \texttt{f} sits
  in front of \texttt{e}, so \texttt{V} pulls back along it; \texttt{g} sits behind \texttt{p} and is total, so it changes which
  value comes out and never which attributes go in, and \texttt{A} passes through.
\end{itemize}

What exactness costs is decidability. \texttt{custom} carries a pair of partial functions, \texttt{held} an
arbitrary \texttt{f}, \texttt{refine} an arbitrary \texttt{q}, so membership of an admissible set is not a question the
declaration can answer. Admission \textbf{specifies} the two families. It does not decide them. What the
discharger row runs is occupancy and option admission, and those are decided.

\textbf{Admission is not what a union consults, and the object is where the two part.} A union admits an
object option by its \textbf{names}: a name set \texttt{N} admits \texttt{object(s,\ ρ)} when every name in \texttt{N} is one
\texttt{s} declares and every \texttt{req}-marked name of \texttt{s} is in \texttt{N}. The union never attempts an object
option's codec. So \texttt{V(c)} and \texttt{A(c)} are \emph{option} admission: admission by names at an object, and
\texttt{dom(e)} and \texttt{dom(p)} at every other term.

Stating the row on a term's own admissible sets instead is unsound and not merely coarse, and one
pair shows it. Take two objects declaring the same single \texttt{req} name and different field terms,
\texttt{⟨a\ :\ String⟩} against \texttt{⟨a\ :\ Number⟩}. Their admissible sets are disjoint on both sides, a map the
first writes carrying a string at \texttt{a} which the second's field term refuses. Both are derivable,
both closed — no free subterm — both lawful, so no other condition declines them. And both options admit
\texttt{⟨a\ ↦\ "1"⟩}, because the union reads the names and never the contents. At an object a term's own
admission is strictly \emph{tighter} than what the union consults, and disjointness of a tighter
predicate implies nothing about disjointness of admission. Off an object the two coincide — being
admitted there \emph{is} being written — so the object is the whole of the difference.

Note what admission by names does \emph{not} read: \texttt{ρ}. An undeclared name disqualifies the option even
where its rest codec would hold it, so an object's open part is unreachable through a union — the
price of deciding by names is that only declared names decide, and an open object that must be an
option should declare what it means to keep. An \texttt{opt} binding \textbf{widens} admission, which makes two
options differing only in optional names the ambiguity to watch: two objects declaring only optional
names both admit the empty map, however disjoint their names. \rulelink{rule:dunion}{\texttt{(dunion)}}'s disjoint-literal premise
is the special case where the discriminator's names decide alone.

\begin{definition}[occupancy]
\label{def:occupancy}
\uses{cond:U}
\lean{Calculus.occV, Calculus.occA}
\leanok
One coarse label per term per side: which constructor of \texttt{Value} the term may write, and
which member of §1 it may store. Occupancy is what decides at the declaration.
\end{definition}

Both over-approximate, and over-approximating is the safe direction — if \texttt{X\ ⊆\ X′}, \texttt{Y\ ⊆\ Y′} and
\texttt{X′\ ∩\ Y′\ =\ ∅} then \texttt{X\ ∩\ Y\ =\ ∅} — so a coarse answer costs a discharge and never buys a wrong one.
Where a term carries an arbitrary function the answer is \texttt{⊤}, every label: \texttt{held}'s \texttt{f} runs before
the subterm and can move a value into any constructor, \texttt{custom} is arbitrary in both directions, and
\texttt{legacy}'s \texttt{r} is arbitrary on the reading side.

\textbf{The writable side's label is the constructor of the value, not the runtime kind}, and the
difference does work. A kind collapses the three set families into one; the constructors do not, and
\texttt{set(string)} writes a set of strings where \texttt{set(number)} writes a set of doubles. Ten constructors
separate what eight kinds cannot. The readable side's label is the member, which the judgement's \texttt{μ}
index already computes.

\texttt{⊤} covers two situations the check treats alike and the semantics does not: \textbf{unknown} — \texttt{held},
\texttt{custom} and \texttt{lazy}, whose answer exists but cannot be read without forcing — and \textbf{everything} —
\texttt{opaque}, which accepts every kind by design. The row separates neither, and both fall to \rulelink{cond:U}{\textbf{(U)}}
per input. The difference resurfaces only inside \texttt{opaque}'s own derivation, where "everything" is
precisely what makes its fixpoint's options pairwise kind-disjoint (§3).

\begin{definition}[the family, conjoined]
\label{def:family}
\uses{def:occupancy,cond:U}
\lean{Calculus.OrthSV, Calculus.OrthSA}
\leanok
An option's set of values is its occupancy intersected with
its option admission, and its set of attributes likewise.
\end{definition}

Disjointness in either factor is enough,
which is what leaves the row decidable where it decides at all: \textbf{on one side, two options separate
when their occupancies on that side are disjoint. On both sides at once, they separate when one
option requires a name the other does not declare.} That is a question about two terms, answered by computation at
the declaration. What it declines is an obligation and not a refusal. No input reaches two
options where option admission is disjoint, and that needs option admission itself, which the held
sets do not determine.

Only \rulelink{cond:D}{\textbf{(D)}} is a premise, and it is the only condition here with a letter, because it appears in
a rule. Everything in the discharger table is a discharger, and the asymmetry is not an
oversight: \rulelink{cond:D}{\textbf{(D)}} is a condition on the judgement's own index, so the rule can demand it, while
a discharger answers a question about inputs that the rule states as \rulelink{cond:U}{\textbf{(U)}} and settles where it
can.
A union no discharger covers is still a term, and what such a union lacks is not derivability but
the theorem: its \rulelink{cond:U}{\textbf{(U)}} is unsettled at the declaration, and the semantics enforces it per input by
failing. Note where the failing lands, and that it is a fact about the two terms rather than a limit
of any table: a container admits its empty case whatever its element admits, so the options collide
at exactly one input — the empty array, the empty map — which is also the one input neither writes
nor reads once they collide, so the loss is that the union declines the empty case and nothing else.
The same union built from \texttt{object}s is decided outright, by names: a \texttt{req} binding is a name the
empty item lacks, so there is no empty case to collide at.

\rulelink{cond:D}{\textbf{(D)}} and the admission row ask the same kind of question of the two sides of the pair, and
neither implies the other. \texttt{union(array(string),\ array(number))} satisfies \rulelink{cond:D}{\textbf{(D)}} — the held sets
are incomparable — and is discharged by nothing: both options write a sequence, both store \texttt{L},
and both admit the empty array.

\texttt{union(set(string),\ set(number))} \emph{is} discharged, and on both sides, because
the writable side is the one that labels by constructor. Those options store \texttt{SS} and \texttt{NS}, so
reading was never in doubt; and they write a set of strings and a set of doubles, which are two
constructors, so the writing is settled as well. What remains is a loss of a different kind: a set's
member is computed from its element, so the options are \texttt{⇝\ SS} and \texttt{⇝\ NS} and the union's own \texttt{μ} is
\texttt{⊤} by the \rulelink{rule:union}{\texttt{(union)}} rule. A \texttt{⊤} term is one only \rulelink{cond:MP}{\textbf{(MP)}} constrains, where a \texttt{⇝\ SS} term is pinned
by \rulelink{cond:W}{\textbf{(W)}} and its read-side corollary. That loss is the union's and not the discharger's.

Two pairs show the two families failing independently, one each way. \texttt{union(set(string),\ set(decimal))} writes a set of strings at both options — a decimal is \emph{held} as a string — and
stores \texttt{SS} at one and \texttt{NS} at the other, so the readable side separates and the writable one does
not. \texttt{union(number,\ decimal)} is the mirror: a number writes a double and a decimal a string, and
both store \texttt{N}, so the writable side separates and the readable one does not. Neither union is
resolved, and each fails on a different half, which is why the two families are asked separately and
never reconciled.

\textbf{Attempting an option answers a real question.} Whether a term \emph{decides its own membership} —
whether "this option accepted the value" is evidence at all — is answered by an admissible set that
\emph{is} \texttt{dom(e)}, and nothing is left for a second condition to add. \texttt{union(literal("A"),\ literal("B"))} is discharged
because \texttt{\{"A"\}} and \texttt{\{"B"\}} do not meet; \texttt{union(literal("A"),\ string)} is not, because \texttt{\{"A"\}} sits
inside every string. A term that decides nothing about its values has the widest domain, and the
widest domain is the one that fails to be disjoint.

Two counterexamples bound the question, and each refutes one form of it:

\begin{itemize}
\tightlist
\item
  \textbf{One side is not enough.} \texttt{union(literal("A"),\ string)} has the left option deciding its own
  membership, so "either option decides" is satisfied, and both options admit \texttt{"A"}. A real answer
  from one option does not resolve a pair; the other option is still answering, and still
  meaninglessly.
\item
  \textbf{Two sides are not enough either.} \texttt{union(array(literal("A")),\ array(literal("B")))} has both
  options as decisive as a term can be, and both admit \texttt{{[}{]}}. Deciding at the element does not
  decide at the container: every container admits its empty case, which no element ever sees.
\end{itemize}

\textbf{A pair sharing a base separates where neither term's own sets do}, which is the discharger
table's fourth row rather than this section's. \texttt{number} and
every numeric literal are the \textbf{same base} under different predicates — the base being every finite
double written as its shortest text, and the predicates \texttt{D} and \texttt{∈\ ℓ}. Two literals over disjoint
value sets are two refinements of one base; refinements of one base read an attribute to the \emph{same}
value, so at most one predicate accepts it, and both sides separate.
\texttt{refine(number,\ positive)} against \texttt{refine(number,\ negative)} is the same argument outside \texttt{literal}.
What carries it is the pair: \texttt{V} and \texttt{A} answer about one term at a time, and this answer is about
two.

\texttt{held} reaches a pair no other row does, and on one side only.
\texttt{union(literal("A"),\ held(literal("B"),\ f,\ g))} separates on the readable side for \emph{every} \texttt{f} and
\texttt{g}, the two sets being \texttt{\{⟨S,\ "A"⟩\}} and \texttt{\{⟨S,\ "B"⟩\}}: the wrapper never decides, the subterm does,
and reading through the wrapper is what lets the question be asked of it. The writable side stays an
obligation, \texttt{f} being free to move a value into any constructor. And
\texttt{union(literal("A"),\ held(literal("A"),\ f,\ g))}, the counterexample below, separates on neither
side. Note that the member disjunct is no help in either case: both options are \texttt{⇝\ S}.

One note on what has no set of its own. \texttt{optional} is a \textbf{field} (§4), not a term,
so nothing passes through it: admission is a function on terms and has no case to give it. And a
union's own admission is not the join of its options' — it is the set where \textbf{exactly one} option
admits, which is the union's own resolution, written as a set. Where two options both admit, the union writes nothing,
which is the empty array above stated the other way round.

\rulelink{cond:D}{\textbf{(D)}} as written is \emph{incomparability}, not disjointness. Two object options \texttt{⟨a\ :\ String⟩} and
\texttt{⟨b\ :\ String⟩} are incomparable yet share the inhabitant \texttt{⟨a\ ↦\ "1",\ b\ ↦\ "2"⟩} — so \rulelink{cond:D}{\textbf{(D)}} admits
them, and resolution falls to the name-set device above. The gap between the two is the set
of unions resolved structurally rather than by type, and \rulelink{cond:U}{\textbf{(U)}} is what fills it.

Stating \rulelink{cond:D}{\textbf{(D)}} as genuine \emph{disjointness} would not be sound either, and the reason is the same
confusion twice: \textbf{held sets are sets of values, and (U)'s readable side quantifies over
attributes.} \texttt{held}
is precisely the term that separates the two — its adapter relabels — so two options can hold
disjoint sets of values and still read one attribute, each to its own value.
\texttt{union(literal("A"),\ held(literal("A"),\ f,\ g))} with \texttt{f} and \texttt{g} swapping \texttt{"B"} for \texttt{"A"} holds
\texttt{\{"A"\}} and \texttt{\{"B"\}}, which are disjoint on the nose, and both options read \texttt{S\ "A"} — the first to
\texttt{"A"}, the second to \texttt{"B"}.

So \rulelink{cond:D}{\textbf{(D)}} is a well-formedness condition on the judgement's \texttt{T} index and not a discharger of
\rulelink{cond:U}{\textbf{(U)}}: no condition on held sets can be one. What decides the readable side is a condition on
\emph{stored attributes} — distinct members \texttt{μᵢ\ ≠\ μⱼ}, or disjoint readable-side admission — which is
what the first two rows of the discharger table take. Attributes rather than values is the point: a
condition on values cannot decide a question about attributes.

\subsection{Soundness of the judgement}\label{theorem-soundness}

\begin{definition}[≼, member agreement]
\label{def:agree}
\uses{cond:MP}
\lean{Calculus.Agree}
\leanok
§2 states \rulelink{cond:MP}{\textbf{(MP)}} per attribute "and recursively"; the proof
needs that as a relation. Define \texttt{a\ ≼\ a′} by recursion on \texttt{a}:
\begin{verbatim}
a ≼ a′  ⟺  member(a) = member(a′)
           ∧  member(a) = L ⟹ pointwise ≼ on shared positions
           ∧  member(a) = M ⟹ pointwise ≼ on shared names
\end{verbatim}
\end{definition}

Set members compare no deeper — an \texttt{SS} carries payloads, not attributes, so there is no member
inside one to disagree on. Shared names, because a reading may drop what a writing then omits: the
discard is the quotient of §2, not a member change. The strong law is \texttt{e(p(a))\ ≼\ a} on \texttt{dom(p)},
and it implies \rulelink{cond:MP}{\textbf{(MP)}} as stated.

\begin{theorem}[soundness]
\label{thm:soundness}
\uses{cond:R,cond:C,cond:W,cond:MP,cond:U,rule:legacy,rule:lazy,rule:union,rule:dunion,rule:opaque}
\lean{Calculus.soundness, Calculus.soundness_core, Calculus.carrier_denote}
\leanok
Every derivable \texttt{⊢\ c\ :\ T\ ⇝\ μ}, and every \texttt{⊢\ f\ :\ T\ ⊗\ π\ ⇝\ μ}, denotes a pair satisfying \rulelink{cond:R}{\textbf{(R)}},
\rulelink{cond:C}{\textbf{(C)}}, \rulelink{cond:W}{\textbf{(W)}}, and \rulelink{cond:MP}{\textbf{(MP)}} in the strong form \texttt{e(p(a))\ ≼\ a}, with three residues,
all named: \rulelink{cond:MP}{\textbf{(MP)}} is off on the attributes a \texttt{legacy} in the derivation rereads; every
\rulelink{rule:union}{\texttt{(union)}} in the derivation carries \rulelink{cond:U}{\textbf{(U)}} \emph{where the union acts} as an obligation, which the
dischargers of this section settle at the declaration wherever they decide a pair, \rulelink{rule:dunion}{\texttt{(dunion)}} and
\texttt{opaque} outright; and a \rulelink{rule:lazy}{\texttt{(lazy)}} whose unfolding re-enters itself \textbf{through a union} is outside
the theorem entirely, for the reason §5 gives under confluent admission and the \rulelink{rule:lazy}{\texttt{(lazy)}} case
below repeats: there is no fixpoint there to prove anything about. Such a union is still derivable.
What it lacks is the theorem, not a rule.
\end{theorem}

The two premises on \rulelink{rule:legacy}{\texttt{(legacy)}}, \texttt{ran(r)\ ⊆\ dom(e)} and
coherence, exist because this proof failed without them. Each excludes a term that the calculus
would otherwise derive and the laws refuse. Null-honesty is not
among them: the null point is an option of \rulelink{rule:union}{\texttt{(union)}}, so what would have been a
premise is that sum's \rulelink{cond:U}{\textbf{(U)}}, and the sugar paragraph of §3 settles it.

\begin{lemma}[injectivity]
\label{lem:injectivity}
\uses{cond:R}
\lean{Calculus.Lawful.e_injective}
\leanok
\rulelink{cond:R}{\textbf{(R)}} makes \texttt{e} injective on \texttt{dom(e)}.
\end{lemma}

\begin{proof}
\uses{cond:R}
\leanok
If \texttt{e(t)\ =\ e(t′)} then \texttt{t\ =\ p(e(t))\ =\ p(e(t′))\ =\ t′}.
\end{proof}

\begin{lemma}[traversal]
\label{lem:traversal}
\uses{cond:R,cond:C}
\lean{Calculus.entrywise, Calculus.lawfulCore_entrywise, Calculus.lawful_entrywise, Calculus.carrier_entrywise}
\leanok
One proof shape serves every product. Let a \emph{step} be an entrywise pair
between the entries of two lists, and lift it entrywise — defined where every entry is, positions
preserved. If the step satisfies a law pointwise, the lift satisfies it listwise.
\end{lemma}

\begin{proof}
\uses{cond:R,cond:C}
\leanok
\rulelink{cond:R}{\textbf{(R)}} because
positions re-read in place, \rulelink{cond:C}{\textbf{(C)}} because each entry re-encodes, \texttt{≼} because a shared position
decodes and re-encodes through the same step, and lengths survive both directions. Keyed products need
no extra clause: \texttt{M} is a name set with a function on it, so keys are unique by type and \texttt{≼}'s \texttt{M}
clause compares at a shared key with no first occurrence to prefer. Proved once; every product case
below is an instantiation.
\end{proof}

\begin{proof}
\leanok
By induction on the derivation, with one refinement at \rulelink{rule:lazy}{\texttt{(lazy)}}. Each case
first says what the pair \emph{is} — the denotation is part of the rule — then discharges the laws. The
grouping is the algebra of §7: points, wrappers, products, sums.

\textbf{Points.} Each ground term writes one member, totally on its held set, and reads exactly that
member back into the held set. \rulelink{cond:W}{\textbf{(W)}} and \rulelink{cond:C}{\textbf{(C)}} are immediate, \texttt{≼} is member equality with no
substructure, and \rulelink{cond:R}{\textbf{(R)}} is the carrier's round trip: identity for \texttt{string}, \texttt{boolean}, \texttt{binary},
\texttt{null} (whose held set is the single value \texttt{null} and whose payload is \texttt{∗}) and \texttt{decimal} (a numeral
in \texttt{Stored} is returned as written, which is what \texttt{Stored} is chosen to make true, so \texttt{Can(decimal)}
is the normalised part of \texttt{N}), shortest-text for \texttt{number} and a numeric \texttt{literal} — the §3 argument that
every finite double survives.

\textbf{Wrappers.}

\emph{held.} \texttt{e\ =\ e\_c\ ∘\ f}, \texttt{p\ =\ g\ ∘\ p\_c}. \rulelink{cond:R}{\textbf{(R)}}: \texttt{g(p\_c(e\_c(f\ t)))\ =\ g(f\ t)\ =\ t}. \rulelink{cond:C}{\textbf{(C)}}:
\texttt{e(p(a))} asks \texttt{f(g(p\_c\ a))\ =\ p\_c\ a\ ∈\ dom(e\_c)}, the premise's \rulelink{cond:C}{\textbf{(C)}}. \texttt{≼}:
\texttt{e(p(a))\ =\ e\_c(p\_c(a))}, the premise's hypothesis. \rulelink{cond:W}{\textbf{(W)}}: the member is untouched.

\emph{refine.} \texttt{e\ =\ e\_c} on \texttt{q}; \texttt{dom(p)\ =\ p\_c⁻¹(q)}. All four laws restrict.

\emph{legacy.} \texttt{e\ =\ e\_c}; \texttt{p\ =\ p\_c\ ∪\ r}, a function \emph{because} of coherence — without that premise the
definition needs a precedence between the two readings, and a precedence is a union with a
privileged write path, the resolution order \rulelink{cond:U}{\textbf{(U)}} exists to refuse. \rulelink{cond:R}{\textbf{(R)}}: \texttt{e\_c(t)\ ∈\ dom(p\_c)}
by the premise's \rulelink{cond:R}{\textbf{(R)}}, and coherence makes any overlap with \texttt{dom(r)} agree, so the composite
reads it to \texttt{t}. \rulelink{cond:C}{\textbf{(C)}}: on \texttt{dom(p\_c)} by the premise; on the extension it \emph{is}
\texttt{ran(r)\ ⊆\ dom(e)} — without it a migrated attribute reads and the write-back is the failure.
\rulelink{cond:W}{\textbf{(W)}}: unchanged. \texttt{≼}: on \texttt{dom(p\_c)} by hypothesis; on the extension it is what the term
is licensed to break, and no further — \texttt{e(r(a))} is canonical, so the breach does not survive its
own round trip, which is the precise sense in which a migrated row behaves like a current one from
the first rewrite on.

\emph{nullable.} Two cases, by the guard of §3, and neither carries a premise.

Where \texttt{owns(c)}, the pair \emph{is} \texttt{c}'s — \texttt{e\ =\ e\_c}, \texttt{p\ =\ p\_c} — and all four laws are the
hypothesis, restated at a \texttt{⊤} index where \rulelink{cond:W}{\textbf{(W)}} is vacuous.

Where it does not, \texttt{e}: \texttt{null\ ↦\ ⟨NULL,\ ∗⟩}, else \texttt{e\_c}; \texttt{p}: \texttt{⟨NULL,\ ∗⟩\ ↦\ null}, else \texttt{p\_c}, and
what were the honesty premises are now available as facts. \rulelink{cond:R}{\textbf{(R)}} at \texttt{null} is immediate; at
\texttt{t\ ≠\ null}, \texttt{e\_c(t)\ ≠\ ⟨NULL,\ ∗⟩}, because if it were, the hypothesis's \rulelink{cond:R}{\textbf{(R)}} would give
\texttt{p\_c(⟨NULL,\ ∗⟩)\ =\ t} and \texttt{owns(c)} would hold. So the \texttt{NULL} branch cannot shadow the reading —
without the guard, \texttt{nullable(held(nullable(string),\ f,\ g))} with \texttt{f} moving a sentinel onto \texttt{null}
would encode the sentinel as \texttt{⟨NULL,\ ∗⟩} and read back \texttt{null}; with it, the inner \texttt{nullable} owns
the point and the outer one adds nothing. \rulelink{cond:C}{\textbf{(C)}}: \texttt{⟨NULL,\ ∗⟩} round-trips; else \texttt{p\_c(a)\ ≠\ null},
because if it were, the hypothesis's \rulelink{cond:C}{\textbf{(C)}} would put \texttt{null} in \texttt{dom(e\_c)} and \texttt{owns(c)} would
hold, so the \texttt{e\_c} branch answers. \texttt{≼} at \texttt{⟨NULL,\ ∗⟩} is exact; elsewhere it is the hypothesis, and
what rules out a non-\texttt{NULL} attribute reading to \texttt{null} and re-encoding as \texttt{⟨NULL,\ ∗⟩} — the member
change \rulelink{cond:MP}{\textbf{(MP)}} exists to refuse — is the same fact.

\texttt{owns(nullable(c))} always, so the second case never nests: §6's idempotence is the construction,
not an equation the proof has to earn.

\emph{custom.} The laws are the premises, verbatim; \rulelink{cond:W}{\textbf{(W)}} is vacuous at \texttt{⊤}.

\emph{lazy.} The pair of \texttt{γ(∗)}. This case is \textbf{three arguments and not one}: what a \texttt{γ}
re-enters itself through decides which one applies.

\textbf{Through a container.} Induct on the pair (input size, derivation size), lexicographically.
Containers strictly shrink the input; the wrappers shrink the derivation and never grow the input.

\textbf{Through wrappers alone.} The measure does not decrease: \texttt{γ(∗)}'s derivation contains \texttt{lazy(γ)}
again, at the same input, so neither component falls and the induction has nothing to say. What
covers this \texttt{γ} is not the induction but partiality — each wrapper moves one way in the information
order, so the unfolding never gains a definedness it did not start with, and it defines nothing at
that input. The laws quantify over \texttt{dom} and hold over an empty one.

\textbf{Through a union.} Outside the theorem, and this is the third residue. Unique admission is a
\emph{negative} condition on the other options, so it is anti-monotone in exactly the argument a
fixpoint must grow: as the recursive options grow, admission grows, and the answer collapses from
\texttt{some} to \texttt{none}. There is no least fixpoint, so there is no pair here to prove the laws of — not a
pair with an empty domain, which is what the wrapper case has, but no semantics at all. That is the
same exclusion §5 records under confluent admission, and it is a residue of this theorem.

Partiality is doing the termination argument's work in the second case, which is why it has to be
its own step.

\textbf{Fields.} A field denotes a pair over presence-extended domains — \texttt{e♭\ :\ T\ ⇀\ 𝔸\ ∪\ \{ε\}},
\texttt{p♭\ :\ (𝔸\ ∪\ \{ε\})\ ⇀\ T}, \texttt{ε} the absent binding — with \texttt{member(ε)} undefined and the laws read over
\texttt{𝔸} as before. \rulelink{rule:field}{\texttt{(field)}} and \rulelink{rule:optional}{\texttt{(optional)}} set \texttt{e♭\ =\ e}, \texttt{p♭\ =\ p}; absence is handled at \rulelink{rule:object}{\texttt{(object)}},
where an \texttt{opt} name absent in the item is a binding absent in the record. A required name
must be present on both current-schema sides.
\texttt{≼} quantifies over \texttt{𝔸}, so \texttt{ε} asks nothing.

\textbf{Products} — each an instantiation of the \rulelink{lem:traversal}{traversal Lemma}.

\emph{array, tuple, tuple-tail, record.} The step is the subterm's pair at each position or name — a
tuple first aligns positions with their codecs, head positionally and tail uniformly, length
mismatches refusing. The Lemma gives all four laws from the induction hypotheses; a \texttt{legacy}
element is where the \rulelink{cond:MP}{\textbf{(MP)}} exemption propagates — pointwise, on the positions it rereads, per
§2.

\emph{object.} The step dispatches each entry by its name: a declared name through its field, with the
\texttt{ε}-handling of the Fields case; an undeclared name by \texttt{ρ}. \texttt{drop} is a filter \emph{before} the step
(§4), so it needs no traversal mode of its own; the presence guards — every \texttt{req} name demanded in
both directions, every non-\texttt{opt} binding demanded of a value — are what keep the step total where
item adds none. \rulelink{cond:W}{\textbf{(W)}}: \texttt{M}. \texttt{≼}: the Lemma's keyed clause; the names \texttt{drop} discards and the
quotient again.

\emph{set.} The premise gives \texttt{e\_c(x)\ =\ ⟨μ,\ w(x)⟩}; the step is \texttt{w} outward and \texttt{p\_c(⟨μ,\ ·⟩)} back, and
the wire set is the image of the held set under it. \rulelink{cond:R}{\textbf{(R)}}: \texttt{w} is injective by the \rulelink{lem:injectivity}{injectivity Lemma} (member fixed,
\texttt{e\_c} injective), so the written set holds one payload per element and the \rulelink{lem:traversal}{traversal Lemma} reads
each back; non-emptiness survives both directions. \rulelink{cond:C}{\textbf{(C)}}: entrywise, and the re-encoded elements
keep their member by \texttt{≼}, so the same payload projection re-succeeds. \texttt{p} need \emph{not} be injective —
two payloads reading to one value shrink the set, which is the quotient again and touches \rulelink{cond:R}{\textbf{(R)}}
nowhere, since canonical payloads are \texttt{w}-images. \rulelink{cond:W}{\textbf{(W)}}: \texttt{μS}, always. \texttt{≼}: member equality, and
a set has no interior members — the container fixes its elements' member, which is why the
element's own \rulelink{cond:MP}{\textbf{(MP)}} is never even invoked inside one.

\textbf{Sums.} Admission, per side: option \texttt{i} admits attribute \texttt{a} when the member/name filter passes
and \texttt{pᵢ(a)} is defined, and admits value \texttt{t} when the shape/name filter passes and \texttt{eᵢ(t)} is
defined. \texttt{p(a)\ =\ pᵢ(a)} for the unique admitting \texttt{i}, else undefined; \texttt{e(t)} dually — \rulelink{cond:U}{\textbf{(U)}} as
semantics rather than aspiration. The laws, and where each needs \rulelink{cond:U}{\textbf{(U)}}:

\begin{itemize}
\tightlist
\item
  \rulelink{cond:R}{\textbf{(R)}}: \texttt{t\ ∈\ dom(e)} names a unique \texttt{i} — that is what makes \texttt{e(t)} defined at all — and \texttt{i}
  admits \texttt{eᵢ(t)} back (the premise's \rulelink{cond:R}{\textbf{(R)}}); \emph{uniqueness} of that admission is \rulelink{cond:U}{\textbf{(U)}} on the
  attribute side, asked at exactly the attributes the union writes. Where it holds,
  \texttt{p(e(t))\ =\ pᵢ(eᵢ(t))\ =\ t}. Partiality cannot hide this: an ambiguous \texttt{e(t)} makes \texttt{p(e(t))}
  undefined, which \rulelink{cond:R}{\textbf{(R)}} counts as failure. What partiality \emph{does} excuse is the input the union
  never wrote: \texttt{t\ ∉\ dom(e)} is outside \rulelink{cond:R}{\textbf{(R)}}'s quantifier, so an ambiguity there costs nothing.
\item
  \rulelink{cond:C}{\textbf{(C)}}: \texttt{a\ ∈\ dom(p)} names a unique \texttt{i}; \texttt{pᵢ(a)\ ∈\ dom(eᵢ)} by the premise's \rulelink{cond:C}{\textbf{(C)}}, so \texttt{i}
  admits it, and uniqueness is \rulelink{cond:U}{\textbf{(U)}} on the value side. Where it holds, \texttt{p(a)\ ∈\ dom(e)}.
\item
  \texttt{≼}: with both uniquenesses the route is the \emph{same option in both directions} — \texttt{i} admits
  \texttt{pᵢ(a)}, so \texttt{e} goes through \texttt{i} — and the premise's hypothesis closes it.
\item
  \rulelink{cond:W}{\textbf{(W)}}: when every \texttt{μᵢ\ =\ μ}, whichever option answers writes \texttt{⟨μ,\ \_⟩}.
\end{itemize}

The dischargers close the residue where they apply:

\emph{dunion.} A canonical item of option \texttt{i} carries the discriminator with a tag in \texttt{Lᵢ}; option
\texttt{j\ ≠\ i} either does not admit the name set or its literal refuses the tag, \texttt{Lⱼ} being disjoint from
\texttt{Lᵢ} — so admission is unique on both sides, and the sum case closes with no residue. The one union
whose \rulelink{cond:U}{\textbf{(U)}} holds by construction, which is what the term is \emph{for}.

\emph{opaque.} The fixpoint of §3: one option per kind, each self-identifying by construction — kinds
partition the values and members partition the attributes — so admission is unique on both sides at
every depth of the fixpoint, and the sum case closes with no residue here either.
\end{proof}

What remains outside the theorem is what §3 already names: \rulelink{rule:custom}{\texttt{(custom)}}'s premises are the
caller's to keep, and each \rulelink{rule:union}{\texttt{(union)}}'s residual \rulelink{cond:U}{\textbf{(U)}}. \texttt{D} is not among them: it is a gate on
both halves of \rulelink{rule:number}{\texttt{(number)}}, and a numeric \texttt{literal}'s base \emph{is} \texttt{number}, so it is inside the theorem.

\begin{center}\rule{0.5\linewidth}{0.5pt}\end{center}

\section{§6 Equations}\label{6-equations}

Which rewrites preserve the pair. Read \texttt{≡} as: the same \texttt{T}, the same \texttt{μ}, and \texttt{e} and \texttt{p} agreeing
everywhere both are defined.

\subsection{Hold}\label{hold}

\begin{theorem}[held(c, id, id) ≡ c]
\label{eq:held-id}
\uses{rule:held}
\lean{Calculus.equiv_through_id, Calculus.equiv_denote_held_id}
\leanok
\texttt{held(c,\ id,\ id)\ ≡\ c}. Identity.
\end{theorem}

\begin{theorem}[isos compose]
\label{eq:held-comp}
\uses{rule:held}
\lean{Calculus.equiv_through_comp, Calculus.equiv_denote_held_comp}
\leanok
\texttt{held(held(c,\ f,\ g),\ f′,\ g′)\ ≡\ held(c,\ f\ ∘\ f′,\ g′\ ∘\ g)}. Isos compose; note the
reversal on the reading half.
\end{theorem}

\begin{theorem}[prisms meet]
\label{eq:refine-meet}
\uses{rule:refine}
\lean{Calculus.equiv_restrict_meet, Calculus.equiv_denote_refine_meet}
\leanok
\texttt{refine(refine(c,\ q),\ q′)\ ≡\ refine(c,\ q\ ∩\ q′)}. Prisms meet.
\end{theorem}

\begin{theorem}[a full refinement is none]
\label{eq:refine-full}
\uses{rule:refine}
\lean{Calculus.equiv_restrict_full, Calculus.equiv_denote_refine_full}
\leanok
\texttt{refine(c,\ T)\ ≡\ c} for \texttt{⊢\ c\ :\ T\ ⇝\ μ}. A refinement admitting everything is no
refinement.
\end{theorem}

\begin{theorem}[a legacy reading that is the current one]
\label{eq:legacy-self}
\uses{rule:legacy}
\lean{Calculus.equiv_extendRead_self, Calculus.equiv_denote_legacy_self}
\leanok
\texttt{legacy(c,\ p\_c)\ ≡\ c}. An older reading that is the current one adds nothing.
\end{theorem}

\begin{theorem}[a migration agrees on the canonical set]
\label{eq:legacy-can}
\uses{rule:legacy,cond:R}
\lean{Calculus.extendRead_eq_on_can, Calculus.denote_legacy_eq_on_can}
\leanok
\texttt{legacy(c,\ r)\ ≈\ c} on \texttt{Can(c)}, for lawful \texttt{c}. A migration differs from a
lawful current encoding only off the canonical set, and the premise it needs is \rulelink{cond:R}{\textbf{(R)}}.
\end{theorem}

\begin{theorem}[choice of one]
\label{eq:union-single}
\uses{rule:union}
\lean{Calculus.equiv_altCodec_single}
\leanok
\texttt{union(c)\ ≡\ c}. Choice of one.
\end{theorem}

\begin{theorem}[union is associative]
\label{eq:union-assoc}
\uses{rule:union}
\lean{Calculus.equiv_altCodec_assoc}
\leanok
\texttt{union(c₁,\ union(c₂,\ c₃))\ ≡\ union(c₁,\ c₂,\ c₃)}. Associative. An ambiguity \emph{inside}
the nested option fails locally on the left and leaves the right side undefined, and \texttt{≡}
compares \texttt{e} and \texttt{p} only where both are defined, so the two sides cannot disagree
there.
\end{theorem}

\begin{theorem}[union is commutative]
\label{eq:union-perm}
\uses{rule:union,cond:U}
\lean{Calculus.equiv_altCodec_perm, Calculus.altCodec_perm, Calculus.equiv_altCodec_swap}
\leanok
\texttt{union(c₁\ …\ cₙ)\ ≡\ union(c\_σ(1)\ …\ c\_σ(n))} for any permutation \texttt{σ}. Commutative:
resolution is unique admission (§5), never declaration order.
\end{theorem}

\begin{theorem}[wrappers preserve the member]
\label{eq:wrapper-mu}
\leanok
preserve the member: a field's member is its attribute's (§4).
\end{theorem}

So wrapping a term never changes the member a key stores.

\subsection{Do not hold}\label{do-not-hold}

\begin{theorem}[a union of arrays is not an array of unions]
\label{eq:union-array}
\uses{rule:union,rule:array,cond:U}
\lean{Calculus.not_equiv_union_array}
\leanok
\texttt{union(array(c₁),\ array(c₂))\ ≢\ array(union(c₁,\ c₂))}. Left: a sequence homogeneous in one
option. Right: a sequence of either, elementwise. Different held sets, and only the right has a
discharger for its \rulelink{cond:U}{\textbf{(U)}}.
\end{theorem}

\begin{theorem}[a tuple with a tail is neither a tuple nor an array]
\label{eq:tuple-tail-neither}
\uses{rule:tuple-tail,rule:tuple,rule:array}
\lean{Calculus.not_equiv_tupleTail_tuple, Calculus.not_equiv_tupleTail_array}
\leanok
\texttt{tuple(c₁\ …\ cₙ,\ tail(c))\ ≢\ tuple(c₁\ …\ cₙ)\ } and \texttt{≢\ array(union(…))}. A fixed head
with a repeating tail is neither: a tuple has no tail and an array has no head.
\end{theorem}

\begin{theorem}[a refinement does not commute with an iso]
\label{eq:refine-held}
\uses{rule:refine,rule:held}
\lean{Calculus.not_equiv_refine_held}
\leanok
\texttt{refine(held(c,\ f,\ g),\ q)\ ≢\ held(refine(c,\ q),\ f,\ g)}. The refinement applies on opposite
sides of the iso, so \texttt{q} is a subset of different sets. Two different terms, both meaningful.
\end{theorem}

\begin{theorem}[a union is not closed under extension]
\label{eq:union-extend}
\uses{rule:union,cond:D,cond:U}
\lean{Calculus.not_equiv_union_extend}
\leanok
\texttt{union(c₁\ …\ cₙ)\ ≢\ union(c₁\ …\ cₙ,\ c)}. Not closed under extension: a further option can
violate \rulelink{cond:D}{\textbf{(D)}} and make the term ill formed, or leave \rulelink{cond:U}{\textbf{(U)}} unresolved where it was
discharged before.
\end{theorem}

\begin{theorem}[lazy is not its own unfolding]
\label{eq:lazy-unfold}
\uses{rule:lazy,rule:set}
\lean{Calculus.not_equiv_lazy}
\leanok
\texttt{lazy(γ)\ ≢\ γ(∗)}. The pair is the same and the μ is not: \texttt{lazy} concludes \texttt{⊤} where
\texttt{γ(∗)} may name its member, so \texttt{set(lazy(γ))} has no derivation. \texttt{≡} compares the index
too.
\end{theorem}

\begin{theorem}[legacy breaks (MP), on purpose]
\label{eq:legacy-mp}
\uses{rule:legacy,cond:MP,cor:readW}
\lean{Calculus.legacy_breaks_MP}
\leanok
\texttt{legacy(c,\ r)} breaks \rulelink{cond:MP}{\textbf{(MP)}}. On purpose, and it is the only term licensed to: changing
what is stored \emph{is} the migration, and §2's read-side corollary falls with it.
\end{theorem}

So a \texttt{legacy} attribute cannot be an index key, which is the price.

\begin{theorem}[the canonizer is not the identity]
\label{eq:canonizer}
\lean{Calculus.e_comp_p_ne_id}
\leanok
\texttt{e\ ∘\ p\ ≢\ id}. Equality holds only on \texttt{Can(c)}.
\end{theorem}

This is §2's quotient.

\begin{tabular}{@{}ll@{}}
\toprule
Equation & Why it is not derivable \\
\midrule
\texttt{set(union(c₁,\ c₂))} — not derivable & The premise on \texttt{set} demands a member, and a union of options writing different members has none. \\
\bottomrule
\end{tabular}

\begin{center}\rule{0.5\linewidth}{0.5pt}\end{center}

\section{§7 The three layers, and the algebra}\label{7-the-three-layers-and-the-algebra}

Three layers answer different questions, and conflating them is what makes a side condition look
optional. A grammar over the terms — a regular tree grammar, with nonterminals indexed by \texttt{μ} — says
which terms can be \emph{written}. It cannot express \rulelink{cond:D}{\textbf{(D)}}, which relates siblings rather than a term to
its subterms, so well-formedness needs the judgement of §3 and §5 rather than a production — nor
the dischargers and \rulelink{cond:U}{\textbf{(U)}}, which are not about the term at all but about the inputs it is given.
And the judgement in turn says nothing about what a term \emph{means}; that is §2, the pair and its laws —
a connection §5's soundness theorem proves rather than leaves implied.

Which places the calculus: \textbf{§1} a universal domain, \textbf{§2} types as quotient retracts of it,
\textbf{§3–§5} a two-index typing relation with side conditions, tied back to §2 by soundness, \textbf{§6} the
equational theory. And the terms are generated by a small algebra, which the soundness proof's
grouping makes official: \textbf{points} (the ground terms); one \textbf{partial iso} (\texttt{held}, with \texttt{refine}
its prism); one \textbf{agreeing read-extension} (\texttt{legacy}); three \textbf{products} — positional (\texttt{array},
\texttt{tuple}) and keyed (\texttt{record}, \texttt{object}), instances of one traversal argument, and unordered (\texttt{set}),
which is the image argument, with injectivity from \rulelink{cond:R}{\textbf{(R)}};
one \textbf{sum} (\texttt{union}), resolved by unique admission, with \rulelink{rule:dunion}{\texttt{(dunion)}} the one whose (U) holds
by construction and \texttt{opaque} the one whose (U) is a theorem; and a \textbf{fixpoint} (\texttt{lazy},
and \texttt{opaque} again). Fields are the keyed product's
presence decoration, and \texttt{nullable} is that sum's guarded insertion (§3) rather than a term of its own.

The algebra makes one more thing visible, which the rules do not show. \textbf{The two halves of the pair are the same construction at two carriers, everywhere the
calculus builds something itself.} A sum filters by the admission \emph{that side} reads and acts
through the option answering alone — the same sentence for \texttt{e} and for \texttt{p}, with value-admission
swapped for attribute-admission. Adjoining the null point is the same construction on values and on
attributes. A product traverses. Reading a discriminator's tag is one lookup, spelled once per
carrier.

Exactly three terms break it, and they are the same three every time: \texttt{held}, \texttt{refine} and
\texttt{legacy} — the wrappers carrying a \textbf{caller's function that is not itself a codec}. \texttt{held}'s \texttt{f}
gates the writing while \texttt{g} only post-processes the reading, which is why the readable side sees
through it and the writable one cannot (§5). \texttt{legacy}'s \texttt{r} widens the reading and touches the writing not at all. \texttt{refine}'s
\texttt{q} is a predicate on \emph{values}, so it gates one carrier directly and the other only after a read.
\texttt{custom} is the degenerate case: both halves are the caller's, so there is no shared construction to
be symmetric about.

So the asymmetry in this calculus is not a property of encoding-versus-decoding. It is the mark of
where a caller's function enters, and it stays local to the four terms that admit one.
That is the divisible /
decidable \cite{divisible-decidable} skeleton, read invariantly — as Invariant rather than as a
contravariant functor — because both directions are present at once.

The judgement settles the laws for every term the calculus builds itself. What it leaves to the
caller is the premise on a function that is not a codec: \rulelink{rule:held}{\texttt{(held)}}'s inverse pair,
\rulelink{rule:refine}{\texttt{(refine)}}'s predicate, \rulelink{rule:legacy}{\texttt{(legacy)}}'s agreement with the reading it
extends, and \rulelink{rule:custom}{\texttt{(custom)}}'s two halves outright. What no derivation settles is each
\texttt{union}'s residual \rulelink{cond:U}{\textbf{(U)}} where no discharger decides it, and how much of any of this an
implementation carries in its own types, which is a question about that implementation and not
about the calculus.

\begin{center}\rule{0.5\linewidth}{0.5pt}\end{center}

\section{Sources}\label{sources}

\begin{thebibliography}{9}
\bibitem{quotient-lenses}
  \textbf{Quotient Lenses.} J. Nathan Foster, Alexandre Pilkiewicz, Benjamin C. Pierce. ICFP 2008. —
  bidirectional transformations well behaved modulo a programmer-controlled equivalence, and
  canonizers.
  \url{https://www.cs.cornell.edu/~jnfoster/papers/quotient-lenses.pdf}

\bibitem{invertible-syntax}
  \textbf{Invertible Syntax Descriptions: Unifying Parsing and Pretty Printing.} Tillmann Rendel, Klaus
  Ostermann. Haskell Symposium 2010. — one description, both directions, over an algebra of partial
  isomorphisms.
  \url{https://www.semanticscholar.org/paper/Invertible-syntax-descriptions:-unifying-parsing-Rendel-Ostermann/44c0cf69d73e4a0673dd04b7da477f7a74f0b735}

\bibitem{ep-pairs}
  \textbf{Embedding-projection pairs.} Retracts of a universal domain, after Scott (1972); standard in
  denotational semantics and in the semantics of gradual types.
  \url{https://www.pls-lab.org/en/Embedding_projection_pair}

\bibitem{divisible-decidable}
  \textbf{Divisible and Decidable.} The contravariant analogues of Applicative and Alternative — products
  and sums on the writing side.
  \url{https://hackage.haskell.org/package/contravariant-1.5.5/docs/Data-Functor-Contravariant-Divisible.html}

\bibitem{tree-transformations}
  \textbf{Combinators for Bidirectional Tree Transformations: A Linguistic Approach to the View-Update
  Problem.} Foster, Greenwald, Moore, Pierce, Schmitt. TOPLAS 29(3), 2007. — the GetPut / PutGet
  laws, and the asymmetric lens the quotient version refines.
  \url{https://dl.acm.org/doi/10.1145/1232420.1232424}

\bibitem{lens-laws}
  \textbf{A Tangled Web of 12 Lens Laws.} A taxonomy of the variants, where the distinctions matter.
  \url{https://link.springer.com/chapter/10.1007/978-3-030-79837-6_11}

\bibitem{dynamodb-types}
  \textbf{Supported data types and naming rules in Amazon DynamoDB.} The \texttt{N} range, stated in exponent
  notation.
  \url{https://docs.aws.amazon.com/amazondynamodb/latest/developerguide/HowItWorks.NamingRulesDataTypes.html}
\end{thebibliography}
